% Complete reviewed Unit 16; editorial master: chapters/ch16/answer-key.md
% Markdown SHA-256: 083eea687565ad2f7d6323b6ee5d86ee7d58bc722284ce995bff11c8a3169eb4
% One-use initial assembly followed by complete content and page review.
\newcommand{\runningtitle}{Unit 16. Answer key}
% Layout inherited from the reviewed previous course; set a running title per unit.
% This is a small hand-maintained TeX layout, not a conversion pipeline.
\RequirePackage{fix-cm}
\documentclass[11pt,twoside]{article}
\usepackage{fontspec}
\setmainfont{Cambria}[
  Renderer=HarfBuzz,
  BoldFont=cambriab.ttf,
  ItalicFont=cambriai.ttf,
  BoldItalicFont=cambriaz.ttf
]
\setmonofont{lmmono10-regular.otf}[BoldFont=lmmonolt10-bold.otf]
\usepackage[paperwidth=165mm,paperheight=241mm,inner=21mm,outer=19mm,
  top=22mm,bottom=20mm,headheight=14pt,headsep=7mm,footskip=10mm]{geometry}
\usepackage{xcolor,array,booktabs,longtable,tabularx}
\usepackage{enumitem,fancyhdr}
\usepackage[unicode,hidelinks]{hyperref}
\definecolor{ink}{HTML}{000000}
\definecolor{muted}{HTML}{4C5558}
\setlength{\parindent}{0pt}
\setlength{\parskip}{3.5pt}
\AtBeginDocument{\fontsize{11}{14.3}\selectfont}
\setlength{\tabcolsep}{4pt}
\renewcommand{\arraystretch}{1.10}
\setlist[itemize]{leftmargin=1.3em,label=\textbullet,itemsep=2pt,topsep=3pt,parsep=0pt}
\setlist[enumerate]{leftmargin=1.8em,itemsep=3pt,topsep=4pt,parsep=0pt}
\widowpenalty=10000
\clubpenalty=10000
\emergencystretch=1em
\raggedbottom
\pagestyle{fancy}
\fancyhf{}
\providecommand{\runningtitle}{Old English}
\fancyhead[LE,RO]{\small\itshape \runningtitle}
\fancyhead[LO,RE]{\small Old English}
\fancyfoot[LE,RO]{\small\thepage}
\renewcommand{\headrulewidth}{0.3pt}
\newcommand{\booktitle}[2]{
  {\fontsize{9}{11}\selectfont\color{muted}#1\par}
  \vspace{3pt}
  {\fontsize{20}{24}\selectfont\color{ink}#2\par}
  \vspace{6pt}
}
\newcommand{\source}[1]{\par{\footnotesize\color{muted}#1\par}}
\newcommand{\readingtext}[1]{
  \par\addvspace{5pt}
  \begingroup\leftskip=1.5em
    \fontsize{12}{16}\selectfont #1\par
  \endgroup\addvspace{5pt}
}
\makeatletter
\renewcommand{\section}{\@startsection{section}{1}{\z@}
  {-2.4ex \@plus -0.6ex \@minus -0.2ex}{0.8ex \@plus 0.2ex}
  {\normalfont\fontsize{12}{15}\selectfont\bfseries\color{ink}}}
\renewcommand{\subsection}{\@startsection{subsection}{2}{\z@}
  {-1.8ex \@plus -0.5ex \@minus -0.2ex}{0.5ex \@plus 0.2ex}
  {\normalfont\fontsize{11}{14}\selectfont\bfseries\color{ink}}}
\makeatother

\hyphenpenalty=10000
\exhyphenpenalty=10000
\emergencystretch=3em
\newcommand{\itempara}[2]{\par\addvspace{3pt}\noindent\hangindent=1.8em\hangafter=1\makebox[1.8em][l]{#1.}#2\par}
\newcommand{\reserve}[1]{\par\penalty-100\begingroup\dimen0=\pagegoal\advance\dimen0 by-\pagetotal\ifdim\dimen0<#1\newpage\fi\endgroup}
\begin{document}
\addfontfeatures{Ligatures=TeXOff}
\raggedright
\booktitle{UNIT 16}{Answer key}

% mdparagraph:P001:begin
The Old English composition models are new teaching material. The reading translation renders Cook’s actual Old English. A model is one defensible answer: preserve the meaning, case relations, agreement, viewpoint and clause connections, while also observing each task’s stated lexical, number and word-order targets. An otherwise valid alternative can miss a requested target.
% mdparagraph:P001:end

% mdparagraph:P002:begin
The exact local helps remain available. A supplied phrase is not a test of unaided recall, even when some of its words have been taught. The form tasks use the selected tables; dictionary-given \textbf{iernan} and \textbf{nearu} are reference inputs, not extra vocabulary-recall tests. Accurate English explanations may differ in wording. N = nominative, G = genitive, D = dative, A = accusative.
% mdparagraph:P002:end

% mdparagraph:P003:begin
Except where the exercise fixes an order, grammatical rearrangements that preserve the relationships are acceptable. Cook also attests such variants as dative \textbf{ðām} beside the selected \textbf{ðǣm}. A spelling variant must be evidenced; it does not excuse a wrong case, person or tense. Distinguish ordinary negative concord from the specially supplied cessation clause in 36.3.
% mdparagraph:P003:end

\reserve{7\baselineskip}
% mdparagraph:P004:begin
\section*{F. Forms and phrases}
% mdparagraph:P004:end

\reserve{8\baselineskip}
% mdparagraph:P005:begin
\subsection*{F01}
% mdparagraph:P005:end

% key: KEY-NEW-U16-F01
% output: NEW-U16-F01-O001
% mdparagraph:P006:begin
\textbf{1. āceorfan: present indicative, you singular.} \textbf{ācierfst} — you cut off (present, singular).
% mdparagraph:P006:end

% output: NEW-U16-F01-O002
% mdparagraph:P007:begin
\textbf{2. āceorfan: preterite indicative, you singular.} \textbf{ācurfe} — you cut off (past, singular).
% mdparagraph:P007:end

% output: NEW-U16-F01-O003
% mdparagraph:P008:begin
\textbf{3. āceorfan: preterite indicative, they.} \textbf{ācurfon} — they cut off (past).
% mdparagraph:P008:end

% output: NEW-U16-F01-O004
% mdparagraph:P009:begin
\textbf{4. āceorfan: explain the stem relations.} The present second singular has \textbf{ie} in \textbf{ācierf-}; the preterite second singular and plural use \textbf{ācurf-}, with \textbf{u}. Add second-singular \textbf{-e} or plural \textbf{-on} to that past stem. Keep the prefix \textbf{ā-} in every form.
% mdparagraph:P009:end

% output: NEW-U16-F01-O005
% mdparagraph:P010:begin
\textbf{5. forlēosan: present indicative, he.} \textbf{forlīest} — he loses.
% mdparagraph:P010:end

% output: NEW-U16-F01-O006
% mdparagraph:P011:begin
\textbf{6. forlēosan: preterite indicative, you singular.} \textbf{forlure} — you lost (singular).
% mdparagraph:P011:end

% output: NEW-U16-F01-O007
% mdparagraph:P012:begin
\textbf{7. forlēosan: preterite indicative, they.} \textbf{forluron} — they lost.
% mdparagraph:P012:end

% output: NEW-U16-F01-O008
% mdparagraph:P013:begin
\textbf{8. forlēosan: explain the stem relations.} The present third singular has \textbf{īe} and the contracted ending \textbf{-st}: \textbf{forlīest}. The preterite second singular and plural have \textbf{u} and \textbf{r}, hence \textbf{forlur-e}, \textbf{forlur-on}. Carrying present \textbf{s} into these past positions misses the stem alternation.
% mdparagraph:P013:end

% output: NEW-U16-F01-O009
% mdparagraph:P014:begin
\textbf{9. ðwēan: present indicative, he.} \textbf{ðwiehð} — he washes.
% mdparagraph:P014:end

% output: NEW-U16-F01-O010
% mdparagraph:P015:begin
\textbf{10. ðwēan: preterite indicative, you singular.} \textbf{ðwōge} — you washed (singular).
% mdparagraph:P015:end

% output: NEW-U16-F01-O011
% mdparagraph:P016:begin
\textbf{11. ðwēan: preterite indicative, they.} \textbf{ðwōgon} — they washed.
% mdparagraph:P016:end

% output: NEW-U16-F01-O012
% mdparagraph:P017:begin
\textbf{12. ðwēan: explain the stem relations.} The present third singular has \textbf{ðwieh-} plus \textbf{-ð}. Preterite second singular and plural have \textbf{ðwōg-}; their endings are \textbf{-e}, \textbf{-on}. The selected first/\allowbreak{}third singular \textbf{ðwōh} has \textbf{h}, so it cannot simply be used as the stem for all past positions.
% mdparagraph:P017:end

% output: NEW-U16-F01-O013
% mdparagraph:P018:begin
\textbf{13. iernan: present indicative, he.} \textbf{iernð} — he runs.
% mdparagraph:P018:end

% output: NEW-U16-F01-O014
% mdparagraph:P019:begin
\textbf{14. iernan: preterite indicative, he.} \textbf{arn} — he ran.
% mdparagraph:P019:end

% output: NEW-U16-F01-O015
% mdparagraph:P020:begin
\textbf{15. iernan: preterite indicative, they.} \textbf{urnon} — they ran.
% mdparagraph:P020:end

% output: NEW-U16-F01-O016
% mdparagraph:P021:begin
\textbf{16. iernan: explain the stem relations.} The selected stems are present \textbf{iern-}, singular preterite \textbf{arn}, and plural preterite \textbf{urn-}. Thus the requested third singular is \textbf{iernð}, the past third singular \textbf{arn}, and the past plural \textbf{urnon}. Neither past stem is formed by merely adding an ending to \textbf{iernan}.
% mdparagraph:P021:end

% mdparagraph:P022:begin
These are the requested indicative positions, not a list of interchangeable principal parts. In particular, \textbf{ācurfe}, \textbf{forlure} and \textbf{ðwōge} are second singular; the corresponding \textbf{-on} forms answer they. The English present/\allowbreak{}past ambiguity of cut does not make the Old English forms interchangeable.
% mdparagraph:P022:end

\reserve{8\baselineskip}
% mdparagraph:P023:begin
\subsection*{F02}
% mdparagraph:P023:end

% key: KEY-NEW-U16-F02
% output: NEW-U16-F02-O001
% mdparagraph:P024:begin
\textbf{1. smierwan: present indicative, I.} \textbf{smierwe} — I anoint.
% mdparagraph:P024:end

% output: NEW-U16-F02-O002
% mdparagraph:P025:begin
\textbf{2. smierwan: present indicative, you singular.} \textbf{smierest} — you anoint (singular).
% mdparagraph:P025:end

% output: NEW-U16-F02-O003
% mdparagraph:P026:begin
\textbf{3. smierwan: preterite indicative, she.} \textbf{smierede} — she anointed.
% mdparagraph:P026:end

% output: NEW-U16-F02-O004
% mdparagraph:P027:begin
\textbf{4. bewęndan: present indicative, I.} \textbf{bewęnde} — I turn.
% mdparagraph:P027:end

% output: NEW-U16-F02-O005
% mdparagraph:P028:begin
\textbf{5. bewęndan: preterite indicative, he.} \textbf{bewęnde} — he turned.
% mdparagraph:P028:end

% output: NEW-U16-F02-O006
% mdparagraph:P029:begin
\textbf{6. plegian: preterite indicative, they.} \textbf{plegodon} — they played.
% mdparagraph:P029:end

% output: NEW-U16-F02-O007
% mdparagraph:P030:begin
\textbf{7. behēafdian: preterite indicative, they.} \textbf{behēafdodon} — they beheaded.
% mdparagraph:P030:end

% output: NEW-U16-F02-O008
% mdparagraph:P031:begin
\textbf{8. gefyllan: uninflected past participle.} \textbf{gefylled} — fulfilled.
% mdparagraph:P031:end

% output: NEW-U16-F02-O009
% mdparagraph:P032:begin
\textbf{9. forgiefan: uninflected past participle.} \textbf{forgiefen} — forgiven.
% mdparagraph:P032:end

% output: NEW-U16-F02-O010
% mdparagraph:P033:begin
\textbf{10. Explain the smierwan stem distribution shown by your answers.} The selected present first singular retains \textbf{w}: \textbf{smierwe}. The present second singular and preterite omit it: \textbf{smierest}, \textbf{smierede}. All three retain the selected \textbf{ie} spelling. The distribution is part of this checked paradigm, not a reason to remove every w from every form.
% mdparagraph:P033:end

% output: NEW-U16-F02-O011
% mdparagraph:P034:begin
\textbf{11. Explain what the two bewęndan answers establish about identifying tense from spelling alone.} \textbf{Bewęnde} represents both present first singular and preterite first/\allowbreak{}third singular. The two requested answers therefore coincide in spelling while differing in tense and person. A subject and context can distinguish them; the ending alone cannot.
% mdparagraph:P034:end

% mdparagraph:P035:begin
\textbf{Gefylled} and \textbf{forgiefen} are past participles, not finite past verbs. Their uninflected use is requested here. In a finite anointing clause, a second-singular past would require \textbf{smieredest}, whereas the requested she form is \textbf{smierede}. Documented i/\allowbreak{}y variants of smierwan are valid spellings but differ from the selected ie series.
% mdparagraph:P035:end

\reserve{8\baselineskip}
% mdparagraph:P036:begin
\subsection*{F03}
% mdparagraph:P036:end

% key: KEY-NEW-U16-F03
% output: NEW-U16-F03-O001
% mdparagraph:P037:begin
\textbf{1. āceorfan: one addressee.} \textbf{āceorf} — cut off, addressed to one.
% mdparagraph:P037:end

% output: NEW-U16-F03-O002
% mdparagraph:P038:begin
\textbf{2. āceorfan: more than one addressee.} \textbf{āceorfað} — cut off, addressed to several.
% mdparagraph:P038:end

% output: NEW-U16-F03-O003
% mdparagraph:P039:begin
\textbf{3. blissian: one addressee.} \textbf{blissa} — rejoice, addressed to one.
% mdparagraph:P039:end

% output: NEW-U16-F03-O004
% mdparagraph:P040:begin
\textbf{4. blissian: more than one addressee.} \textbf{blissiað} — rejoice, addressed to several.
% mdparagraph:P040:end

% output: NEW-U16-F03-O005
% mdparagraph:P041:begin
\textbf{5. smierwan: one addressee.} \textbf{smiere} — anoint, addressed to one.
% mdparagraph:P041:end

% output: NEW-U16-F03-O006
% mdparagraph:P042:begin
\textbf{6. smierwan: more than one addressee.} \textbf{smierwað} — anoint, addressed to several.
% mdparagraph:P042:end

% output: NEW-U16-F03-O007
% mdparagraph:P043:begin
\textbf{7. ārīsan: one addressee.} \textbf{ārīs} — arise, addressed to one.
% mdparagraph:P043:end

% output: NEW-U16-F03-O008
% mdparagraph:P044:begin
\textbf{8. ārīsan: more than one addressee.} \textbf{ārīsað} — arise, addressed to several.
% mdparagraph:P044:end

% output: NEW-U16-F03-O009
% mdparagraph:P045:begin
\textbf{9. forlǣtan: one addressee.} \textbf{forlǣt} — leave, addressed to one.
% mdparagraph:P045:end

% output: NEW-U16-F03-O010
% mdparagraph:P046:begin
\textbf{10. forlǣtan: more than one addressee.} \textbf{forlǣtað} — leave, addressed to several.
% mdparagraph:P046:end

% output: NEW-U16-F03-O011
% mdparagraph:P047:begin
\textbf{11. geswīcan: one addressee.} \textbf{geswīc} — cease, addressed to one.
% mdparagraph:P047:end

% output: NEW-U16-F03-O012
% mdparagraph:P048:begin
\textbf{12. geswīcan: more than one addressee.} \textbf{geswīcað} — cease, addressed to several.
% mdparagraph:P048:end

% output: NEW-U16-F03-O013
% mdparagraph:P049:begin
\textbf{13. ęfthweorfan: one addressee.} \textbf{ęfthweorf} — return, addressed to one.
% mdparagraph:P049:end

% output: NEW-U16-F03-O014
% mdparagraph:P050:begin
\textbf{14. ęfthweorfan: more than one addressee.} \textbf{ęfthweorfað} — return, addressed to several.
% mdparagraph:P050:end

% mdparagraph:P051:begin
The singular commands do not all have the same ending: compare \textbf{blissa}, \textbf{smiere} and \textbf{ārīs}. The plural follows the selected present stem, including \textbf{smierwað} with w. These are ordinary command pairs; the separately taught shortened plural before following gē is not requested here.
% mdparagraph:P051:end

\reserve{8\baselineskip}
% mdparagraph:P052:begin
\subsection*{F04}
% mdparagraph:P052:end

% key: KEY-NEW-U16-F04
% output: NEW-U16-F04-O001
% mdparagraph:P053:begin
\textbf{1. findan: present subjunctive, third singular.} \textbf{finde} — that he find.
% mdparagraph:P053:end

% output: NEW-U16-F04-O002
% mdparagraph:P054:begin
\textbf{2. findan: present subjunctive, third plural.} \textbf{finden} — that they find.
% mdparagraph:P054:end

% output: NEW-U16-F04-O003
% mdparagraph:P055:begin
\textbf{3. fēdan: present subjunctive, third singular.} \textbf{fēde} — that he feed.
% mdparagraph:P055:end

% output: NEW-U16-F04-O004
% mdparagraph:P056:begin
\textbf{4. fēdan: present subjunctive, third plural.} \textbf{fēden} — that they feed.
% mdparagraph:P056:end

% output: NEW-U16-F04-O005
% mdparagraph:P057:begin
\textbf{5. iernan: plain infinitive.} \textbf{iernan} — to run, plain form.
% mdparagraph:P057:end

% output: NEW-U16-F04-O006
% mdparagraph:P058:begin
\textbf{6. iernan: inflected infinitive after tō (give the infinitive only).} \textbf{iernanne} — to run, inflected form; give no separate tō here.
% mdparagraph:P058:end

% mdparagraph:P059:begin
Present subjunctive singular \textbf{-e} contrasts with plural \textbf{-en}. Do not substitute indicative third-singular or plural endings. The infinitives have \textbf{-an} and \textbf{-anne}; neither is a present participle. In the last answer tō is already specified by the question, so the requested form is \textbf{iernanne} alone.
% mdparagraph:P059:end

\reserve{8\baselineskip}
% mdparagraph:P060:begin
\subsection*{F05}
% mdparagraph:P060:end

% key: KEY-NEW-U16-F05
% output: NEW-U16-F05-O001
% mdparagraph:P061:begin
\textbf{1. nearu: masculine nominative singular.} \textbf{nearu} — narrow, N masculine singular.
% mdparagraph:P061:end

% output: NEW-U16-F05-O002
% mdparagraph:P062:begin
\textbf{2. nearu: masculine accusative singular.} \textbf{nearone} — narrow, A masculine singular.
% mdparagraph:P062:end

% output: NEW-U16-F05-O003
% mdparagraph:P063:begin
\textbf{3. nearu: masculine genitive singular.} \textbf{nearwes} — narrow, G masculine singular.
% mdparagraph:P063:end

% output: NEW-U16-F05-O004
% mdparagraph:P064:begin
\textbf{4. nearu: masculine dative plural.} \textbf{nearwum} — narrow, D masculine plural.
% mdparagraph:P064:end

% output: NEW-U16-F05-O005
% mdparagraph:P065:begin
\textbf{5. both sons, accusative.} \textbf{bēgen suna} — both sons, A plural.
% mdparagraph:P065:end

% output: NEW-U16-F05-O006
% mdparagraph:P066:begin
\textbf{6. of both sons, genitive.} \textbf{bēgra suna} — of both sons, G plural.
% mdparagraph:P066:end

% output: NEW-U16-F05-O007
% mdparagraph:P067:begin
\textbf{7. to both sons, dative.} \textbf{bǣm sunum} — to both sons, D plural.
% mdparagraph:P067:end

% output: NEW-U16-F05-O008
% mdparagraph:P068:begin
\textbf{8. Explain the nearu stem alternation shown by your answers.} The masculine nominative singular ends in \textbf{-u}. The selected accusative has \textbf{-o-} before \textbf{-ne}: \textbf{nearone}. Before \textbf{-es} and \textbf{-um} the stem appears as \textbf{nearw-}: \textbf{nearwes}, \textbf{nearwum}. Do not build all four forms from an unchanged nearu-.
% mdparagraph:P068:end

% output: NEW-U16-F05-O009
% mdparagraph:P069:begin
\textbf{9. Explain the grammatical number of both in the phrases.} \textbf{Bēgen} has plural forms, although its meaning refers to two. \textbf{Bēgen}, \textbf{bēgra}, \textbf{bǣm} agree with the plural forms of \textbf{sunu}. It has no separate dual declension; with a dual personal pronoun such as \textbf{unc}, the appropriate form of both is still plural.
% mdparagraph:P069:end

% mdparagraph:P070:begin
The three phrases require no article. \textbf{Suna} is syncretic: in the first phrase it is accusative plural, in the second genitive plural. \textbf{Bēgen} versus \textbf{bēgra} makes that case contrast visible. \textbf{Bǣm sunum} has dative plural agreement in both words.
% mdparagraph:P070:end

\reserve{7\baselineskip}
% mdparagraph:P071:begin
\section*{A. Reference, clause structure and time}
% mdparagraph:P071:end

\reserve{8\baselineskip}
% mdparagraph:P072:begin
\subsection*{A01}
% mdparagraph:P072:end

% key: KEY-NEW-U16-A01
% output: NEW-U16-A01-O001
% mdparagraph:P073:begin
\textbf{1. First reading.} \textbf{The king went.} Hē can refer back to the king, the subject of the first sentence.
% mdparagraph:P073:end

% output: NEW-U16-A01-O002
% mdparagraph:P074:begin
\textbf{2. Second reading.} \textbf{The son went.} Hē can instead refer to the king’s son. Both candidates are masculine singular; proximity to the pronoun alone does not settle the reference.
% mdparagraph:P074:end

% output: NEW-U16-A01-O003
% mdparagraph:P075:begin
\textbf{3. First explicit version.} \textbf{Se cyning ēode.} Only the pronoun has been replaced; the action and tense remain unchanged.
% mdparagraph:P075:end

% output: NEW-U16-A01-O004
% mdparagraph:P076:begin
\textbf{4. Second explicit version.} \textbf{Se sunu ēode.} The explicit nominative subject selects the son. The two readings and their corresponding rewrites may be presented in the opposite order.
% mdparagraph:P076:end

% mdparagraph:P077:begin
The explicit noun phrases must be nominative subjects. A nearest-noun rule would erase the real ambiguity; replacing more than hē would also miss the restricted rewriting task.
% mdparagraph:P077:end

\reserve{8\baselineskip}
% mdparagraph:P078:begin
\subsection*{A02}
% mdparagraph:P078:end

% key: KEY-NEW-U16-A02
% output: NEW-U16-A02-O001
% mdparagraph:P079:begin
\textbf{1. Sentence 1 translation.} “The king saw the son go.” “Saw the son walking” also conveys the perceived action; it need not introduce an Old English progressive construction.
% mdparagraph:P079:end

% output: NEW-U16-A02-O002
% mdparagraph:P080:begin
\textbf{2. Sentence 2 translation.} “The king knew that the son went.” “Had gone” is defensible if an anterior context is stated, but it is not forced by the spelling ēode alone.
% mdparagraph:P080:end

% output: NEW-U16-A02-O003
% mdparagraph:P081:begin
\textbf{3. Sentence 1 actor.} \textbf{Ðone sunu} is accusative masculine singular. It denotes the actor of the plain infinitive \textbf{gān} in the perception construction: the son, not the king, does the going.
% mdparagraph:P081:end

% output: NEW-U16-A02-O004
% mdparagraph:P082:begin
\textbf{4. Sentence 2 actor.} \textbf{Se sunu} is nominative masculine singular, the subject of finite \textbf{ēode} inside the \textbf{ðæt} content clause. The subject’s case is determined within that clause.
% mdparagraph:P082:end

% output: NEW-U16-A02-O005
% mdparagraph:P083:begin
\textbf{5. Verb-form comparison.} \textbf{Gān} is a nonfinite plain infinitive. \textbf{Ēode} is a finite third-singular preterite of gān here; indicative and subjunctive share that spelling. A factual indicative or contextual report-subjunctive analysis is possible, provided the answer acknowledges that the surface alone does not distinguish them.
% mdparagraph:P083:end

\reserve{8\baselineskip}
% mdparagraph:P084:begin
\subsection*{A03}
% mdparagraph:P084:end

% key: KEY-NEW-U16-A03
% output: NEW-U16-A03-O001
% mdparagraph:P085:begin
\textbf{1. Sentence 1 translation.} “He said that the king was there.” The king’s presence is simultaneous with the saying, as the specified context requires.
% mdparagraph:P085:end

% output: NEW-U16-A03-O002
% mdparagraph:P086:begin
\textbf{2. Sentence 1 finite forms.} \textbf{Sægde:} sęcgan, third-singular preterite indicative. \textbf{Wǣre:} wesan, third-singular preterite subjunctive in the reported clause. The subject se cyning rules out the homographic second-singular indicative analysis of wǣre.
% mdparagraph:P086:end

% output: NEW-U16-A03-O003
% mdparagraph:P087:begin
\textbf{3. Sentence 1 nonfinite forms.} There is \textbf{no infinitive or participle}. Both sægde and wǣre are finite.
% mdparagraph:P087:end

% output: NEW-U16-A03-O004
% mdparagraph:P088:begin
\textbf{4. Sentence 1 construction and time.} The content is a finite \textbf{ðæt} clause with reported \textbf{wǣre}. Its past form is interpreted from the past saying viewpoint; the context places the presence at that same time. It does not mean that the king had already ceased to be there. English backshift alone cannot determine the relative time.
% mdparagraph:P088:end

% output: NEW-U16-A03-O005
% mdparagraph:P089:begin
\textbf{5. Sentence 2 translation.} “He said that the king had made the ship.” The making is complete before the saying viewpoint.
% mdparagraph:P089:end

% output: NEW-U16-A03-O006
% mdparagraph:P090:begin
\textbf{6. Sentence 2 finite forms.} \textbf{Sægde:} third-singular preterite indicative of sęcgan. \textbf{Hæfde:} third-singular preterite of habban, finite auxiliary in the content clause. The model selects the report subjunctive; its indicative spelling is identical, so a factual-indicative interpretation is also defensible with that limit explained.
% mdparagraph:P090:end

% output: NEW-U16-A03-O007
% mdparagraph:P091:begin
\textbf{7. Sentence 2 nonfinite form.} \textbf{Geworht} is the uninflected past participle of \textbf{gewyrcean}. It combines with finite \textbf{hæfde} and accusative patient \textbf{ðæt scip}. It is not itself a finite verb.
% mdparagraph:P091:end

% output: NEW-U16-A03-O008
% mdparagraph:P092:begin
\textbf{8. Sentence 2 construction and time.} The transitive \textbf{habban + past participle} group presents the making as complete at the past viewpoint. That complete group, together with the stated context, supports anterior “had made.” Calling both verbs simply past would miss the relation between them; this is not a passive with the ship as subject.
% mdparagraph:P092:end

% output: NEW-U16-A03-O009
% mdparagraph:P093:begin
\textbf{9. Sentence 3 translation.} “He said that the king intended to sing.” “Would sing” is acceptable if it retains the intended later action rather than asserting its eventual fulfillment.
% mdparagraph:P093:end

% output: NEW-U16-A03-O010
% mdparagraph:P094:begin
\textbf{10. Sentence 3 finite forms.} \textbf{Sægde:} third-singular preterite indicative of sęcgan. \textbf{Wolde:} third-singular preterite of willan, finite volitional verb in the content clause. The selected report-subjunctive analysis shares its spelling with the indicative; a contextual indicative report is possible if that syncretism is acknowledged.
% mdparagraph:P094:end

% output: NEW-U16-A03-O011
% mdparagraph:P095:begin
\textbf{11. Sentence 3 nonfinite form.} \textbf{Singan} is the plain infinitive complement of \textbf{wolde}. Its understood actor is the king.
% mdparagraph:P095:end

% output: NEW-U16-A03-O012
% mdparagraph:P096:begin
\textbf{12. Sentence 3 construction and time.} \textbf{Willan + plain infinitive} expresses volition or intention. In the stated context, singing is still intended and lies after the saying. Neither wolde nor the English word would proves that the singing happened. This contrasts with simultaneous wǣre in sentence 1 and completed geworht hæfde in sentence 2.
% mdparagraph:P096:end

\reserve{7\baselineskip}
% mdparagraph:P097:begin
\section*{C. Short compositions}
% mdparagraph:P097:end

\reserve{8\baselineskip}
% mdparagraph:P098:begin
\subsection*{C01}
% mdparagraph:P098:end

% key: KEY-NEW-U16-C01
% output: NEW-U16-C01-O001
% mdparagraph:P099:begin
\textbf{1. Speech and command.} \textbf{Se cyning sægde his suna: ‘Sęle mē ðā bōc.’}
% mdparagraph:P099:end

% mdparagraph:P100:begin
The king said to his son, “Give me the book.” \textbf{His suna} is the dative speech recipient; \textbf{sęle} is a singular command. \textbf{Mē} is the recipient of giving and \textbf{ðā bōc} its accusative object.
% mdparagraph:P100:end

% output: NEW-U16-C01-O002
% mdparagraph:P101:begin
\textbf{2. Executed action.} \textbf{Se sunu sealde him ðā bōc.}
% mdparagraph:P101:end

% mdparagraph:P102:begin
The son gave the book to the king. \textbf{Sealde} narrates the completed action; \textbf{him} remains dative. An explicit \textbf{ðǣm cyninge} would remove reference ambiguity, but would replace the recipient pronoun that the next part asks you to explain.
% mdparagraph:P102:end

% output: NEW-U16-C01-O003
% mdparagraph:P103:begin
\textbf{3. Recipient reference.} Quoted \textbf{mē} and narrative \textbf{him} both denote the king. The first is first person because the king is speaking; the second is third person because the narrator is describing him. Changing the pronoun’s person need not change its referent.
% mdparagraph:P103:end

\reserve{8\baselineskip}
% mdparagraph:P104:begin
\subsection*{C02}
% mdparagraph:P104:end

% key: KEY-NEW-U16-C02
% output: NEW-U16-C02-O001
% mdparagraph:P105:begin
\textbf{1. Statement 1.} \textbf{Ic cann ðone mann.} — “I know that man.” The personal object is accusative, and cunnan means acquaintance here. Another evidenced knowing verb could be valid language, but the prompt specifically requests cunnan.
% mdparagraph:P105:end

% output: NEW-U16-C02-O002
% mdparagraph:P106:begin
\textbf{2. Statement 2.} \textbf{Ic cann singan.} — “I know how to sing.” The plain infinitive supplies the learned skill; no separate translation of English how is required in this frame.
% mdparagraph:P106:end

% output: NEW-U16-C02-O003
% mdparagraph:P107:begin
\textbf{3. Statement 3.} \textbf{Ic wāt ðæt hē hēr is.} — “I know that he is here.” Witan takes a finite content clause; \textbf{hē} is its nominative subject. The model presents the content factually with indicative is. A subjunctive needs an appropriate contextual interpretation, not an automatic rule that every dependent clause must have one.
% mdparagraph:P107:end

% output: NEW-U16-C02-O004
% mdparagraph:P108:begin
\textbf{4. Complement comparison.} The first complement is an \textbf{accusative person}, the second a \textbf{plain infinitive}, and the third a \textbf{finite ðæt clause}. These examples contrast particular senses and constructions. They do not establish a universal persons-versus-facts division: witan can also take a personal object, as in Cook 35.7.
% mdparagraph:P108:end

\reserve{8\baselineskip}
% mdparagraph:P109:begin
\subsection*{C03}
% mdparagraph:P109:end

% key: KEY-NEW-U16-C03
% output: NEW-U16-C03-O001
% mdparagraph:P110:begin
\textbf{1. Whole comparison.} \textbf{Ðā męnn plegodon wel; se cyning plegode bęt ðonne hīe dydon.}
% mdparagraph:P110:end

% mdparagraph:P111:begin
“The men played well; the king played better than they did.” \textbf{Plegodon} and \textbf{dydon} are plural; \textbf{plegode} is singular. \textbf{Bęt} compares the playing, and ðonne introduces its comparison clause.
% mdparagraph:P111:end

% output: NEW-U16-C03-O002
% mdparagraph:P112:begin
\textbf{2. Recovered action and agreement.} \textbf{Dydon} stands for \textbf{played} and agrees with \textbf{hīe}, the men. Its subject is not the singular king. Repeating \textbf{plegodon} could make a grammatical comparison, but would miss the expressly requested substitute-dōn target. This use does not create modern English do-support.
% mdparagraph:P112:end

\clearpage
% mdparagraph:P113:begin
\section*{Cook XII, XIII.32}
% mdparagraph:P113:end

\par\begin{minipage}{\linewidth}\raggedright
% mdparagraph:P114:begin
\subsection*{Cook XII, XIII.32.1}
% mdparagraph:P114:end

% key: KEY-EX-XII-XIII-32-01
% mdparagraph:P115:begin
Ic grēte ēow ealle.
% mdparagraph:P115:end

% output: EX-XII-XIII-32-01-O001
% mdparagraph:P116:begin
\textbf{Meaning:} I bid all of you farewell.
% mdparagraph:P116:end

% mdparagraph:P117:begin
\textbf{Grēte} uses the farewell sense of grētan with accusative \textbf{ēow ealle}. Ealle agrees with the persons addressed; it is not an adverb meaning entirely. This direct farewell meets the named lexical target without adding a separate bidding construction.
% mdparagraph:P117:end

\end{minipage}\par\addvspace{4pt}

\par\begin{minipage}{\linewidth}\raggedright
% mdparagraph:P118:begin
\subsection*{Cook XII, XIII.32.2}
% mdparagraph:P118:end

% key: KEY-EX-XII-XIII-32-02
% mdparagraph:P119:begin
Ðā hīe rēowon swilce twēntig furlanga oððe ðrītig, ðā gesāwon hīe Hǣlend uppan ðǣm sǣ gān.
% mdparagraph:P119:end

% output: EX-XII-XIII-32-02-O001
% mdparagraph:P120:begin
\textbf{Meaning:} When they had rowed about twenty or thirty furlongs, they saw Jesus walking on the sea.
% mdparagraph:P120:end

% mdparagraph:P121:begin
\textbf{Rēowon} is a simple preterite with anterior meaning supplied by the sequence; English had does not require an Old English perfect here. \textbf{Hǣlend} is the accusative actor of \textbf{gān}, and the infinitive stays last as directed. The supplied approximation and quantity phrase retains both twenty and thirty; \textbf{uppan ðǣm sǣ} uses dative agreement with the selected masculine sea noun. “Walking” in English does not make gān a participle.
% mdparagraph:P121:end

\end{minipage}\par\addvspace{4pt}

\par\begin{minipage}{\linewidth}\raggedright
% mdparagraph:P122:begin
\subsection*{Cook XII, XIII.32.3}
% mdparagraph:P122:end

% key: KEY-EX-XII-XIII-32-03
% mdparagraph:P123:begin
Hē ārās and bebēad ðǣm winde, and cwæð tō ðǣm sǣ, Geswīc.
% mdparagraph:P123:end

% output: EX-XII-XIII-32-03-O001
% mdparagraph:P124:begin
\textbf{Meaning:} He arose, commanded the wind, and said to the sea, “Cease.”
% mdparagraph:P124:end

% mdparagraph:P125:begin
All three actions remain: \textbf{ārās}, \textbf{bebēad}, \textbf{cwæð}. The wind is the dative recipient \textbf{ðǣm winde}, while the sea follows \textbf{tō}. \textbf{Geswīc} is the singular absolute command to cease; it needs no genitive complement here. Substituting an accusative wind or treating the command as a past statement changes the structure.
% mdparagraph:P125:end

\end{minipage}\par\addvspace{4pt}

\par\begin{minipage}{\linewidth}\raggedright
% mdparagraph:P126:begin
\subsection*{Cook XII, XIII.32.4}
% mdparagraph:P126:end

% key: KEY-EX-XII-XIII-32-04
% mdparagraph:P127:begin
Se wind geswāc, and wearð miclu smyltnes.
% mdparagraph:P127:end

% output: EX-XII-XIII-32-04-O001
% mdparagraph:P128:begin
\textbf{Meaning:} The wind ceased, and a great calm arose.
% mdparagraph:P128:end

% mdparagraph:P129:begin
\textbf{Geswāc} is absolute cessation. \textbf{Wearð} means came about, with nominative \textbf{miclu smyltnes}; there is no participle and therefore no passive auxiliary construction. \textbf{Miclu} is the selected feminine nominative strong form. The English indefinite article is omitted as instructed.
% mdparagraph:P129:end

\end{minipage}\par\addvspace{4pt}

\par\begin{minipage}{\linewidth}\raggedright
% mdparagraph:P130:begin
\subsection*{Cook XII, XIII.32.5}
% mdparagraph:P130:end

% key: KEY-EX-XII-XIII-32-05
% mdparagraph:P131:begin
Ðū hīerst his stefne, and ðū ne wāst hwanon hē cymð, ne hwider hē gǣð.
% mdparagraph:P131:end

% output: EX-XII-XIII-32-05-O001
% mdparagraph:P132:begin
\textbf{Meaning:} You hear his voice and do not know where he comes from or where he goes.
% mdparagraph:P132:end

% mdparagraph:P133:begin
\textbf{Wāst} is the selected second-singular present of witan; separate \textbf{ne} negates it. The two dependent questions retain \textbf{hwanon … cymð} and \textbf{hwider … gǣð}. Use the uncontracted form supplied by the taught system; an invented negative contraction does not improve compliance. The voice \textbf{stefne} is accusative object of hear.
% mdparagraph:P133:end

\end{minipage}\par\addvspace{4pt}

\par\begin{minipage}{\linewidth}\raggedright
% mdparagraph:P134:begin
\subsection*{Cook XII, XIII.32.6}
% mdparagraph:P134:end

% key: KEY-EX-XII-XIII-32-06
% mdparagraph:P135:begin
Hē geseah twēgen brōðor hiera nętt on ðone sǣ sęndan; hīe wǣron fisceras.
% mdparagraph:P135:end

% output: EX-XII-XIII-32-06-O001
% mdparagraph:P136:begin
\textbf{Meaning:} He saw two brothers casting their net into the sea; they were fishermen.
% mdparagraph:P136:end

% mdparagraph:P137:begin
\textbf{Twēgen brōðor} is the accusative actor phrase of perception \textbf{sęndan}. Their \textbf{nętt} is the singular object sent; \textbf{on ðone sǣ} gives the goal. The finite \textbf{wǣron} belongs to the second clause. A finite sending clause could express related content, but it would replace the selected perception-infinitive construction.
% mdparagraph:P137:end

\end{minipage}\par\addvspace{4pt}

\par\begin{minipage}{\linewidth}\raggedright
% mdparagraph:P138:begin
\subsection*{Cook XII, XIII.32.7}
% mdparagraph:P138:end

% key: KEY-EX-XII-XIII-32-07
% mdparagraph:P139:begin
And hē sægde him, Cumað æfter mē, and ic dō ðæt git bēoð manna fisceras; and hīe sōna forlēton hiera nętt.
% mdparagraph:P139:end

% output: EX-XII-XIII-32-07-O001
% mdparagraph:P140:begin
\textbf{Meaning:} He told them, “Come after me, and I will make the two of you fishers of men”; they immediately left their nets.
% mdparagraph:P140:end

% mdparagraph:P141:begin
\textbf{Him} is the dative recipient of the saying. The result is a finite \textbf{ðæt} clause: nominative dual \textbf{git} takes plural \textbf{bēoð}, and \textbf{fisceras} is its nominative predicate. Do not replace git with an accusative object of dōn; this is the taught make-that construction. The model omits a will word as required, retains genitive \textbf{manna}, and uses plural \textbf{nętt} without an invented -u ending. The following-person phrase and immediately remain explicit.
% mdparagraph:P141:end

\end{minipage}\par\addvspace{4pt}

\par\begin{minipage}{\linewidth}\raggedright
% mdparagraph:P142:begin
\subsection*{Cook XII, XIII.32.8}
% mdparagraph:P142:end

% key: KEY-EX-XII-XIII-32-08
% mdparagraph:P143:begin
Gif ðīn hand oððe ðīn fōt ðē ǣswicað, āceorf hit of, and weorp fram ðē; bętere ðē is ðæt ðū healt on līf gā, ðonne ðū twā handa and twēgen fēt hæbbe and on ēce fȳr sīe worpen.
% mdparagraph:P143:end

% output: EX-XII-XIII-32-08-O001
% mdparagraph:P144:begin
\textbf{Meaning:} If your hand or foot causes you to stumble, cut it off and cast it away; it is better for you to enter life lame than to have two hands and two feet and be cast into eternal fire.
% mdparagraph:P144:end

% mdparagraph:P145:begin
Singular \textbf{ǣswicað} agrees with the alternatives hand or foot; \textbf{ðē} is the offended person, accusative. \textbf{Hit} resumes the alternative body part, and \textbf{of} is the off particle. The second command omits it as directed. In the supplied better-clause, \textbf{ðē} is dative; keep \textbf{ðæt} before the first complement, omit it after \textbf{ðonne}, and retain \textbf{gā}, \textbf{hæbbe}, \textbf{sīe worpen}. The final words are a passive; worpen is the selected prior participle.
% mdparagraph:P145:end

\end{minipage}\par\addvspace{4pt}

\par\begin{minipage}{\linewidth}\raggedright
% mdparagraph:P146:begin
\subsection*{Cook XII, XIII.32.9}
% mdparagraph:P146:end

% key: KEY-EX-XII-XIII-32-09
% mdparagraph:P147:begin
Gif hwilc mann hæfð hund scēapa and forlīest ān of ðǣm, ne forlǣtt hē ðā nigon and hundnigontig on muntum and gǣð and sēcð ðæt ðe forwearð?
% mdparagraph:P147:end

% output: EX-XII-XIII-32-09-O001
% mdparagraph:P148:begin
\textbf{Meaning:} If a man has a hundred sheep and loses one, does he not leave the ninety-nine in the mountains and go to seek the one that was perishing?
% mdparagraph:P148:end

% mdparagraph:P149:begin
Keep the five present actions distinct: \textbf{hæfð}, \textbf{forlīest}, \textbf{forlǣtt}, \textbf{gǣð}, \textbf{sēcð}. The final \textbf{forwearð} is a preterite in the identifying relative; the English was perishing does not require a progressive auxiliary. \textbf{Hund scēapa} has the supplied quantity construction, and \textbf{ān of ðǣm} selects one from the set. The question keeps \textbf{ne} but has no equivalent of auxiliary does.
% mdparagraph:P149:end

\end{minipage}\par\addvspace{4pt}

\par\begin{minipage}{\linewidth}\raggedright
% mdparagraph:P150:begin
\subsection*{Cook XII, XIII.32.10}
% mdparagraph:P150:end

% key: KEY-EX-XII-XIII-32-10
% mdparagraph:P151:begin
And gif hit gelimpð ðæt hē hit finde, sōðlice ic ēow sęcge ðæt hē swīðor ofer ðæt blissað ðonne ofer ðā nigon and hundnigontig.
% mdparagraph:P151:end

% output: EX-XII-XIII-32-10-O001
% mdparagraph:P152:begin
\textbf{Meaning:} If he happens to find it, truly I tell you that he rejoices more over that sheep than over the ninety-nine.
% mdparagraph:P152:end

% mdparagraph:P153:begin
The first \textbf{hit} in the supplied \textbf{gif hit gelimpð} is impersonal; the second is the sheep. \textbf{Finde} is present subjunctive, not indicative findað. \textbf{Swīðor} compares the degree of rejoicing; both objects retain \textbf{ofer} with accusative. Replacing the comparison with a statement of preference would lose this use of more.
% mdparagraph:P153:end

\end{minipage}\par\addvspace{4pt}

\par\begin{minipage}{\linewidth}\raggedright
% mdparagraph:P154:begin
\subsection*{Cook XII, XIII.32.11}
% mdparagraph:P154:end

% key: KEY-EX-XII-XIII-32-11
% mdparagraph:P155:begin
Ðonne hē hām cymð, hē cleopað his frīend and his nēahgebūras tōgædere and cwið, Blissiað mid mē, for-ðon-ðe ic mīn forloren scēap fand.
% mdparagraph:P155:end

% output: EX-XII-XIII-32-11-O001
% mdparagraph:P156:begin
\textbf{Meaning:} When he comes home, he calls his friends and neighbours together and says, “Rejoice with me, because I found my lost sheep.”
% mdparagraph:P156:end

% mdparagraph:P157:begin
\textbf{Hām} expresses homeward motion without an added preposition. Singular \textbf{cleopað} summons the friends; plural \textbf{Blissiað} addresses them. The quoted first-singular preterite is the course-selected \textbf{fand}, omitting have as directed. \textbf{Forloren} is strong accusative neuter singular with \textbf{mīn … scēap}; the participial base does not license an arbitrary -e ending. The documented first-singular funde is a linguistic variant, but differs from the selected cell.
% mdparagraph:P157:end

\end{minipage}\par\addvspace{4pt}

\clearpage
% mdparagraph:P158:begin
\section*{Cook XII, XIII.33}
% mdparagraph:P158:end

\par\begin{minipage}{\linewidth}\raggedright
% mdparagraph:P159:begin
\subsection*{Cook XII, XIII.33.1}
% mdparagraph:P159:end

% key: KEY-EX-XII-XIII-33-01
% mdparagraph:P160:begin
And, intō ðǣm hūse gangende, hīe fundon ðæt bearn mid Marian his mēder.
% mdparagraph:P160:end

% output: EX-XII-XIII-33-01-O001
% mdparagraph:P161:begin
\textbf{Meaning:} Going into the house, they found the child with Mary, his mother.
% mdparagraph:P161:end

% mdparagraph:P162:begin
The whole \textbf{intō ðǣm hūse gangende} phrase is supplied; gangende accompanies the plural actors, while \textbf{fundon} is finite. \textbf{Bearn} is the accusative child. \textbf{Mid Marian his mēder} retains the supplied name form and the dative relation, with his mother identifying Mary. Do not turn the going into an unrelated extra finite subject.
% mdparagraph:P162:end

\end{minipage}\par\addvspace{4pt}

\par\begin{minipage}{\linewidth}\raggedright
% mdparagraph:P163:begin
\subsection*{Cook XII, XIII.33.2}
% mdparagraph:P163:end

% key: KEY-EX-XII-XIII-33-02
% mdparagraph:P164:begin
Ārīs, and nim ðæt bearn and his mōdor on niht, and flēoh tō Egypta lande, and bēo ðǣr oð ðæt ic ðē sęcge.
% mdparagraph:P164:end

% output: EX-XII-XIII-33-02-O001
% mdparagraph:P165:begin
\textbf{Meaning:} Arise, take the child and his mother by night, flee to the land of the Egyptians, and remain there until I tell you.
% mdparagraph:P165:end

% mdparagraph:P166:begin
Preserve all the commands and both people taken. \textbf{Ārīs}, \textbf{nim}, \textbf{flēoh}, \textbf{bēo} address one person. The supplied \textbf{on niht} gives the time, while \textbf{tō Egypta lande} gives the destination. \textbf{Sęcge} is the finite subjunctive in the until clause and must stand last there; it is not another command to the hearer.
% mdparagraph:P166:end

\end{minipage}\par\addvspace{4pt}

\par\begin{minipage}{\linewidth}\raggedright
% mdparagraph:P167:begin
\subsection*{Cook XII, XIII.33.3}
% mdparagraph:P167:end

% key: KEY-EX-XII-XIII-33-03
% mdparagraph:P168:begin
Hē ðā ārās and nōm ðæt bearn and his mōdor on niht and fōr on Egypta land, ðæt gefylled wǣre ðæt fram dryhtne ðurh ðone wītgan cweden wæs, Of Egyptum ic mīnne sunu cleopode.
% mdparagraph:P168:end

% output: EX-XII-XIII-33-03-O001
% mdparagraph:P169:begin
\textbf{Meaning:} He arose, took the child and his mother by night, and journeyed into the land of the Egyptians, so that what the Lord spoke through the prophet might be fulfilled: “From the Egyptians I called my son.”
% mdparagraph:P169:end

% mdparagraph:P170:begin
The selected narrative forms are \textbf{ārās}, \textbf{nōm}, \textbf{fōr}. \textbf{On Egypta land} is a goal, unlike the time expression on niht. The purpose has \textbf{gefylled wǣre}, with no separate might word; the thing fulfilled is the spoken utterance. \textbf{Fram dryhtne} and \textbf{ðurh ðone wītgan} retain their distinct cases. The quoted \textbf{cleopode} is first singular without have. A separately evidenced travel verb could express journeyed, but no new fēran system is required by this unqualified action.
% mdparagraph:P170:end

\end{minipage}\par\addvspace{4pt}

\par\begin{minipage}{\linewidth}\raggedright
% mdparagraph:P171:begin
\subsection*{Cook XII, XIII.33.4}
% mdparagraph:P171:end

% key: KEY-EX-XII-XIII-33-04
% mdparagraph:P172:begin
Hwæs magon hīe biddan?
% mdparagraph:P172:end

% output: EX-XII-XIII-33-04-O001
% mdparagraph:P173:begin
\textbf{Meaning:} What can they ask for?
% mdparagraph:P173:end

% mdparagraph:P174:begin
\textbf{Hwæs} is genitive: the thing requested is governed by \textbf{biddan}. \textbf{Magon} is plural and governs the plain infinitive. Do not substitute accusative \textbf{hwæt} for the genitive thing requested merely because English what has no visible case ending. The same request frame can also take an accusative person, as in 36.8.
% mdparagraph:P174:end

\end{minipage}\par\addvspace{4pt}

\par\begin{minipage}{\linewidth}\raggedright
% mdparagraph:P175:begin
\subsection*{Cook XII, XIII.33.5}
% mdparagraph:P175:end

% key: KEY-EX-XII-XIII-33-05
% mdparagraph:P176:begin
Sēo sunne gīemð ðæs dæges.
% mdparagraph:P176:end

% output: EX-XII-XIII-33-05-O001
% mdparagraph:P177:begin
\textbf{Meaning:} The sun rules over the day.
% mdparagraph:P177:end

% mdparagraph:P178:begin
\textbf{Gīemð} is singular, and its rule/\allowbreak{}govern sense takes genitive \textbf{ðæs dæges}. This is lexical government, not a general rule that everything following an English over must be genitive. The already taught gīeman frame meets Cook’s §156f target.
% mdparagraph:P178:end

\end{minipage}\par\addvspace{4pt}

\par\begin{minipage}{\linewidth}\raggedright
% mdparagraph:P179:begin
\subsection*{Cook XII, XIII.33.6}
% mdparagraph:P179:end

% key: KEY-EX-XII-XIII-33-06
% mdparagraph:P180:begin
Hīe ne wēndon nāhtes būtan dēaðes selfes.
% mdparagraph:P180:end

% output: EX-XII-XIII-33-06-O001
% mdparagraph:P181:begin
\textbf{Meaning:} They expected nothing except death itself.
% mdparagraph:P181:end

% mdparagraph:P182:begin
\textbf{Wēndon} means expected, with a genitive thing expected: \textbf{nāhtes}. The supplied exceptive phrase keeps \textbf{dēaðes selfes} in the corresponding genitive relation. \textbf{Ne} participates in the negative statement. A finite that-clause is another construction of wēnan, but does not explain the genitives required here.
% mdparagraph:P182:end

\end{minipage}\par\addvspace{4pt}

\par\begin{minipage}{\linewidth}\raggedright
% mdparagraph:P183:begin
\subsection*{Cook XII, XIII.33.7}
% mdparagraph:P183:end

% key: KEY-EX-XII-XIII-33-07
% mdparagraph:P184:begin
Ðū geswīcst ðīnes weorces.
% mdparagraph:P184:end

% output: EX-XII-XIII-33-07-O001
% mdparagraph:P185:begin
\textbf{Meaning:} You will cease from your work.
% mdparagraph:P185:end

% mdparagraph:P186:begin
\textbf{Geswīcst} is present second singular with the future meaning supplied by context. \textbf{Ðīnes weorces} is the genitive complement of cease. A willan paraphrase may express futurity elsewhere, but the selected simple present already fulfils the task; adding an English-shaped from phrase would miss the requested genitive government.
% mdparagraph:P186:end

\end{minipage}\par\addvspace{4pt}

\par\begin{minipage}{\linewidth}\raggedright
% mdparagraph:P187:begin
\subsection*{Cook XII, XIII.33.8}
% mdparagraph:P187:end

% key: KEY-EX-XII-XIII-33-08
% mdparagraph:P188:begin
Wēn is ðæt hē frēond finde.
% mdparagraph:P188:end

% output: EX-XII-XIII-33-08-O001
% mdparagraph:P189:begin
\textbf{Meaning:} Perhaps he will find a friend.
% mdparagraph:P189:end

% mdparagraph:P190:begin
The taught \textbf{Wēn is ðæt} pattern expresses expectation or possibility, followed by present subjunctive \textbf{finde}. \textbf{Frēond} is the accusative object found. The English will is not translated automatically by willan, which would introduce an intention not required by this model.
% mdparagraph:P190:end

\end{minipage}\par\addvspace{4pt}

\par\begin{minipage}{\linewidth}\raggedright
% mdparagraph:P191:begin
\subsection*{Cook XII, XIII.33.9}
% mdparagraph:P191:end

% key: KEY-EX-XII-XIII-33-09
% mdparagraph:P192:begin
Hē sægde him, Ic wille gān on fiscnoð; hīe sægdon, And wē willað gān mid ðē.
% mdparagraph:P192:end

% output: EX-XII-XIII-33-09-O001
% mdparagraph:P193:begin
\textbf{Meaning:} He said to them, “I will go fishing.” They replied, “And we will go with you.”
% mdparagraph:P193:end

% mdparagraph:P194:begin
The quotation changes from \textbf{ic wille} to \textbf{wē willað}, retaining each speaker’s own viewpoint. The supplied activity phrase is \textbf{on fiscnoð}; the companions follow \textbf{mid} in \textbf{mid ðē}. Do not leave the second willan singular merely because the first quotation was singular.
% mdparagraph:P194:end

\end{minipage}\par\addvspace{4pt}

\par\begin{minipage}{\linewidth}\raggedright
% mdparagraph:P195:begin
\subsection*{Cook XII, XIII.33.10}
% mdparagraph:P195:end

% key: KEY-EX-XII-XIII-33-10
% mdparagraph:P196:begin
Hīe ēodon ūt and ēodon on scip and ne fēngon nāwiht on ðǣre niht.
% mdparagraph:P196:end

% output: EX-XII-XIII-33-10-O001
% mdparagraph:P197:begin
\textbf{Meaning:} They went out, went into a ship, and caught nothing that night.
% mdparagraph:P197:end

% mdparagraph:P198:begin
Both movements and the failed catch remain finite past events. \textbf{On scip} is the goal, whereas \textbf{on ðǣre niht} is the supplied time expression. \textbf{Ne … nāwiht} gives one negative meaning, not two cancelling negatives. The uninflected-looking \textbf{niht} is the selected dative here; do not invent nihte to make its case more visible.
% mdparagraph:P198:end

\end{minipage}\par\addvspace{4pt}

\clearpage
% mdparagraph:P199:begin
\section*{Cook XII, XIII.34}
% mdparagraph:P199:end

\par\begin{minipage}{\linewidth}\raggedright
% mdparagraph:P200:begin
\subsection*{Cook XII, XIII.34.1}
% mdparagraph:P200:end

% key: KEY-EX-XII-XIII-34-01
% mdparagraph:P201:begin
Swā hwā swā ðū sīe, ne forlǣt mē.
% mdparagraph:P201:end

% output: EX-XII-XIII-34-01-O001
% mdparagraph:P202:begin
\textbf{Meaning:} Whoever you are, do not forsake me.
% mdparagraph:P202:end

% mdparagraph:P203:begin
\textbf{Swā hwā swā} introduces free choice with subjunctive \textbf{sīe}. \textbf{Ne forlǣt} is a singular negative command and takes accusative \textbf{mē}. The -e of sīe does not make it an infinitive, and the negative command needs no auxiliary equivalent to English do.
% mdparagraph:P203:end

\end{minipage}\par\addvspace{4pt}

\par\begin{minipage}{\linewidth}\raggedright
% mdparagraph:P204:begin
\subsection*{Cook XII, XIII.34.2}
% mdparagraph:P204:end

% key: KEY-EX-XII-XIII-34-02
% mdparagraph:P205:begin
Gif ðū nānne ne finde ðe ðē fēde, ęfthweorf hider.
% mdparagraph:P205:end

% output: EX-XII-XIII-34-02-O001
% mdparagraph:P206:begin
\textbf{Meaning:} If you find nobody who will feed you, return here.
% mdparagraph:P206:end

% mdparagraph:P207:begin
\textbf{Nānne} is the possible person found, with negative \textbf{ne}; hypothetical \textbf{finde} is present subjunctive. The relative \textbf{ðe ðē fēde} describes a possible provider, so fēde is also subjunctive, with the fed person accusative. Singular \textbf{ęfthweorf} is the command to return. The two subjunctives express this context; they are not a universal translation of English will.
% mdparagraph:P207:end

\end{minipage}\par\addvspace{4pt}

\par\begin{minipage}{\linewidth}\raggedright
% mdparagraph:P208:begin
\subsection*{Cook XII, XIII.34.3}
% mdparagraph:P208:end

% key: KEY-EX-XII-XIII-34-03
% mdparagraph:P209:begin
Genyhtsumien mīna bēc unc bǣm.
% mdparagraph:P209:end

% output: EX-XII-XIII-34-03-O001
% mdparagraph:P210:begin
\textbf{Meaning:} Let my books suffice for both of us.
% mdparagraph:P210:end

% mdparagraph:P211:begin
The complete \textbf{Genyhtsumien … unc bǣm} construction is supplied. \textbf{Mīna bēc} is nominative plural, so the optative/\allowbreak{}subjunctive is plural \textbf{genyhtsumien}. \textbf{Unc} is dual dative; \textbf{bǣm} is dative plural because both has no separate dual table. Cook’s verified noun is books; the cited reading has possessions and provides the construction, not a replacement for the exercise noun.
% mdparagraph:P211:end

\end{minipage}\par\addvspace{4pt}

\par\begin{minipage}{\linewidth}\raggedright
% mdparagraph:P212:begin
\subsection*{Cook XII, XIII.34.4}
% mdparagraph:P212:end

% key: KEY-EX-XII-XIII-34-04
% mdparagraph:P213:begin
Bēgen ðā cnapan sindon wel gelǣrede.
% mdparagraph:P213:end

% output: EX-XII-XIII-34-04-O001
% mdparagraph:P214:begin
\textbf{Meaning:} Both boys are well educated.
% mdparagraph:P214:end

% mdparagraph:P215:begin
\textbf{Bēgen ðā cnapan} and predicate \textbf{gelǣrede} agree in the masculine plural. Here gelǣred is the taught lexical adjective educated or trained. Do not infer a separately assigned finite gelǣran conjugation from its participial origin. Well trained is an appropriate contextual rendering of the same adjective.
% mdparagraph:P215:end

\end{minipage}\par\addvspace{4pt}

\par\begin{minipage}{\linewidth}\raggedright
% mdparagraph:P216:begin
\subsection*{Cook XII, XIII.34.5}
% mdparagraph:P216:end

% key: KEY-EX-XII-XIII-34-05
% mdparagraph:P217:begin
Hīe wǣron bēgen rihtwīse beforan Gode.
% mdparagraph:P217:end

% output: EX-XII-XIII-34-05-O001
% mdparagraph:P218:begin
\textbf{Meaning:} They were both righteous before God.
% mdparagraph:P218:end

% mdparagraph:P219:begin
\textbf{Bēgen} and \textbf{rihtwīse} are plural predicates referring to hīe; \textbf{wǣron} supplies the past tense. \textbf{Beforan Gode} has a dative complement. Translating both by a dual personal pronoun would change they into we two or you two rather than supplying the required agreeing word.
% mdparagraph:P219:end

\end{minipage}\par\addvspace{4pt}

\par\begin{minipage}{\linewidth}\raggedright
% mdparagraph:P220:begin
\subsection*{Cook XII, XIII.34.6}
% mdparagraph:P220:end

% key: KEY-EX-XII-XIII-34-06
% mdparagraph:P221:begin
Ic hālsie ðē ðurh ðone libbendan God ðæt ðū ūs sęcge hwæðer ðū sīe Crīst, Godes Sunu.
% mdparagraph:P221:end

% output: EX-XII-XIII-34-06-O001
% mdparagraph:P222:begin
\textbf{Meaning:} I solemnly charge you by the living God to tell us whether you are Christ, the Son of God.
% mdparagraph:P222:end

% mdparagraph:P223:begin
\textbf{Hālsie} takes accusative \textbf{ðē} and the oath phrase \textbf{ðurh ðone libbendan God}. That participial adjective is supplied as a weak accusative masculine form. \textbf{Ūs} is the dative recipient of telling; \textbf{sęcge} belongs to the request clause. \textbf{Hwæðer} introduces the dependent whether-question, with \textbf{sīe}. Do not confuse the charged person’s accusative with the speech recipient’s dative.
% mdparagraph:P223:end

\end{minipage}\par\addvspace{4pt}

\par\begin{minipage}{\linewidth}\raggedright
% mdparagraph:P224:begin
\subsection*{Cook XII, XIII.34.7}
% mdparagraph:P224:end

% key: KEY-EX-XII-XIII-34-07
% mdparagraph:P225:begin
Ic ēow sęcge, Æfter ðissum gē gesēoð Mannes Sunu on ðā swīðran hand Godes mægenðrymmes sittan and on heofones wolcnum cuman.
% mdparagraph:P225:end

% output: EX-XII-XIII-34-07-O001
% mdparagraph:P226:begin
\textbf{Meaning:} I tell you: after these things you will see the Son of Man sitting at the right hand of God’s majesty and coming in the clouds of heaven.
% mdparagraph:P226:end

% mdparagraph:P227:begin
One accusative actor, \textbf{Mannes Sunu}, controls both perception infinitives \textbf{sittan} and \textbf{cuman}. English sitting/\allowbreak{}coming may therefore be rendered without participles. The supplied right-hand phrase retains Cook’s expressly required accusative \textbf{on ðā swīðran hand}; do not change it to a motion meaning just to explain the case. \textbf{On heofones wolcnum} is a separate locative phrase. The model omits the specified things and articles.
% mdparagraph:P227:end

\end{minipage}\par\addvspace{4pt}

\par\begin{minipage}{\linewidth}\raggedright
% mdparagraph:P228:begin
\subsection*{Cook XII, XIII.34.8}
% mdparagraph:P228:end

% key: KEY-EX-XII-XIII-34-08
% mdparagraph:P229:begin
Sīe inc bǣm æfter incrum gelēafan.
% mdparagraph:P229:end

% output: EX-XII-XIII-34-08-O001
% mdparagraph:P230:begin
\textbf{Meaning:} May it be for you two according to your faith.
% mdparagraph:P230:end

% mdparagraph:P231:begin
\textbf{Sīe} expresses the wish without inserting the omitted it. \textbf{Inc} is dative dual, while \textbf{bǣm} agrees in its available dative plural form. \textbf{Incrum} preserves the dual possessive reference. The supplied \textbf{æfter … gelēafan} means according to, rather than following someone in space.
% mdparagraph:P231:end

\end{minipage}\par\addvspace{4pt}

\par\begin{minipage}{\linewidth}\raggedright
% mdparagraph:P232:begin
\subsection*{Cook XII, XIII.34.9}
% mdparagraph:P232:end

% key: KEY-EX-XII-XIII-34-09
% mdparagraph:P233:begin
Sōðlice Mannes Sunu is tō cumanne on his Fæder wuldre, and ðonne āgylt hē ǣlcum be his āgnum weorce.
% mdparagraph:P233:end

% output: EX-XII-XIII-34-09-O001
% mdparagraph:P234:begin
\textbf{Meaning:} Truly the Son of Man is to come in his Father’s glory and will then reward everyone according to his own work.
% mdparagraph:P234:end

% mdparagraph:P235:begin
Both \textbf{is tō cumanne} and the complete reward clause are visibly supplied. The former is this bounded be-to construction, not merely the previously taught purpose after a coming verb. \textbf{Cumanne} is the selected form; manuscript evidence also has cumenne. \textbf{C06:} Cook’s printed hint \textbf{āgyltan,113} is corrected to \textbf{āgyldan}, strong III, for reward/\allowbreak{}repay. Its contracted third singular \textbf{āgylt} is supported, and \textbf{ǣlcum} is the dative recipient. Keep the original hint identifiable, but do not manufacture a weak reward form from it. No shall word is added.
% mdparagraph:P235:end

\end{minipage}\par\addvspace{4pt}

\par\begin{minipage}{\linewidth}\raggedright
% mdparagraph:P236:begin
\subsection*{Cook XII, XIII.34.10}
% mdparagraph:P236:end

% key: KEY-EX-XII-XIII-34-10
% mdparagraph:P237:begin
Ðonne Crīst cymð, ðonne ne wāt nān mann hwanon hē is.
% mdparagraph:P237:end

% output: EX-XII-XIII-34-10-O001
% mdparagraph:P238:begin
\textbf{Meaning:} When Christ comes, no one then knows where he is from.
% mdparagraph:P238:end

% mdparagraph:P239:begin
The temporal pair is \textbf{ðonne … ðonne}. \textbf{Ne wāt nān mann} is a single negative statement; wāt is the selected third-singular present, with a dependent whence-question. A future interpretation of comes does not require a separate future auxiliary or an unassigned contracted negative-witan form.
% mdparagraph:P239:end

\end{minipage}\par\addvspace{4pt}

\par\begin{minipage}{\linewidth}\raggedright
% mdparagraph:P240:begin
\subsection*{Cook XII, XIII.34.11}
% mdparagraph:P240:end

% key: KEY-EX-XII-XIII-34-11
% mdparagraph:P241:begin
Hǣlend cleopode and lǣrde on ðǣm temple and cwæð, Mē gē cunnon, and gē witon hwanon ic eom; and ic ne cōm of mē selfum, ac sē is sōð ðe mē sęnde, ðone gē ne cunnon.
% mdparagraph:P241:end

% output: EX-XII-XIII-34-11-O001
% mdparagraph:P242:begin
\textbf{Meaning:} Jesus cried out and taught in the temple: “You know me and where I am from. I did not come from myself; the one who sent me is true, and you do not know him.”
% mdparagraph:P242:end

% mdparagraph:P243:begin
The model retains crying, teaching and saying. \textbf{Cunnan} has an accusative person; \textbf{witan} has the whence-content. In \textbf{sē … ðe mē sęnde}, the sender is nominative subject; in \textbf{ðone gē ne cunnon}, the same person is accusative object. \textbf{Of mē selfum} keeps the dative self-phrase after of. The supplied predicate \textbf{sōð} describes the sender as true; it is not an adverb modifying sent.
% mdparagraph:P243:end

\end{minipage}\par\addvspace{4pt}

\par\begin{minipage}{\linewidth}\raggedright
% mdparagraph:P244:begin
\subsection*{Cook XII, XIII.34.12}
% mdparagraph:P244:end

% key: KEY-EX-XII-XIII-34-12
% mdparagraph:P245:begin
Gē hine ne cūðon; ic hine cann, and gif ic sęcge ðæt ic hine ne cunne, ic bēo ēow gelīc; ac ic hine cann, and ic healde his sprǣce.
% mdparagraph:P245:end

% output: EX-XII-XIII-34-12-O001
% mdparagraph:P246:begin
\textbf{Meaning:} You have not known him; I know him. If I say that I do not know him, I will be like you; but I know him and keep his saying.
% mdparagraph:P246:end

% mdparagraph:P247:begin
The first clause uses the requested preterite \textbf{cūðon} without have. Reported \textbf{cunne} is subjunctive and final in its clause; the surrounding \textbf{cann} forms are indicative. \textbf{Bēo} expresses the future without shall, and \textbf{ēow} is dative with like. \textbf{Healde his sprǣce} uses keep/\allowbreak{}observe with an accusative saying, rather than the earlier physical hold sense.
% mdparagraph:P247:end

\end{minipage}\par\addvspace{4pt}

\par\begin{minipage}{\linewidth}\raggedright
% mdparagraph:P248:begin
\subsection*{Cook XII, XIII.34.13}
% mdparagraph:P248:end

% key: KEY-EX-XII-XIII-34-13
% mdparagraph:P249:begin
Gif ic mē selfne wuldrie, nis mīn wuldor nāwiht; mīn Fæder is sē ðe mē wuldrað, be ðǣm gē cweðað ðæt hē ēower God sīe.
% mdparagraph:P249:end

% output: EX-XII-XIII-34-13-O001
% mdparagraph:P250:begin
\textbf{Meaning:} If I glorify myself, my glory is nothing; my Father is the one who glorifies me, and you say concerning him that he is your God.
% mdparagraph:P250:end

% mdparagraph:P251:begin
\textbf{Mē selfne} is accusative reflexive. The following negative clause begins with \textbf{nis}, satisfying the verb-first control, and \textbf{nāwiht} is its supplied predicate. \textbf{Be ðǣm} means concerning whom. The final reported \textbf{sīe} is subjunctive. Cook allows a g-containing weak-II spelling series, but the selected first singular is \textbf{wuldrie}, not an unnoticed mixture of spellings.
% mdparagraph:P251:end

\end{minipage}\par\addvspace{4pt}

\clearpage
% mdparagraph:P252:begin
\section*{Cook XII, XIII.35}
% mdparagraph:P252:end

\par\begin{minipage}{\linewidth}\raggedright
% mdparagraph:P253:begin
\subsection*{Cook XII, XIII.35.1}
% mdparagraph:P253:end

% key: KEY-EX-XII-XIII-35-01
% mdparagraph:P254:begin
Fōh ðisne ðoðor.
% mdparagraph:P254:end

% output: EX-XII-XIII-35-01-O001
% mdparagraph:P255:begin
\textbf{Meaning:} Catch this ball.
% mdparagraph:P255:end

% mdparagraph:P256:begin
\textbf{Fōh} is the singular command. \textbf{Ðisne ðoðor} is accusative masculine singular; the demonstrative agrees with the ball’s grammatical gender. The source’s rare ball word is a noun, not a finite verb inferred from its ending.
% mdparagraph:P256:end

\end{minipage}\par\addvspace{4pt}

\par\begin{minipage}{\linewidth}\raggedright
% mdparagraph:P257:begin
\subsection*{Cook XII, XIII.35.2}
% mdparagraph:P257:end

% key: KEY-EX-XII-XIII-35-02
% mdparagraph:P258:begin
Hwā fēng ðone blacan ðoðor?
% mdparagraph:P258:end

% output: EX-XII-XIII-35-02-O001
% mdparagraph:P259:begin
\textbf{Meaning:} Who caught the black ball?
% mdparagraph:P259:end

% mdparagraph:P260:begin
\textbf{Hwā} is the subject and \textbf{fēng} the third-singular preterite. The object \textbf{ðone blacan ðoðor} has weak accusative masculine adjective agreement after the demonstrative. Present catch and preterite caught require different fōn stems; the question does not change that distinction.
% mdparagraph:P260:end

\end{minipage}\par\addvspace{4pt}

\par\begin{minipage}{\linewidth}\raggedright
% mdparagraph:P261:begin
\subsection*{Cook XII, XIII.35.3}
% mdparagraph:P261:end

% key: KEY-EX-XII-XIII-35-03
% mdparagraph:P262:begin
Gā ðæs weges ðe ðū cōme.
% mdparagraph:P262:end

% output: EX-XII-XIII-35-03-O001
% mdparagraph:P263:begin
\textbf{Meaning:} Go along the way by which you came.
% mdparagraph:P263:end

% mdparagraph:P264:begin
The route is genitive \textbf{ðæs weges}, the pattern directly encountered in the reading. \textbf{Gā} is the singular command; \textbf{cōme} is preterite indicative second singular in the relative clause. Its final -e alone does not make it subjunctive. The selected genitive target should not be replaced by an unverified accusative route merely because English has no preposition.
% mdparagraph:P264:end

\end{minipage}\par\addvspace{4pt}

\par\begin{minipage}{\linewidth}\raggedright
% mdparagraph:P265:begin
\subsection*{Cook XII, XIII.35.4}
% mdparagraph:P265:end

% key: KEY-EX-XII-XIII-35-04
% mdparagraph:P266:begin
Ic ne plegode nǣfre bęt ðonne ic tō-dæg dyde.
% mdparagraph:P266:end

% output: EX-XII-XIII-35-04-O001
% mdparagraph:P267:begin
\textbf{Meaning:} I never played better than I did today.
% mdparagraph:P267:end

% mdparagraph:P268:begin
\textbf{Plegode} expresses ordinary playing. \textbf{Dyde} substitutes for that same action and has first-singular agreement. \textbf{Bęt … ðonne} compares actions, with \textbf{tō-dæg} at the current speaker’s viewpoint. A repeated playing verb could form a valid comparison, but the selected model practises substitute dōn. \textbf{Ne … nǣfre} is one negative assertion.
% mdparagraph:P268:end

\end{minipage}\par\addvspace{4pt}

\par\begin{minipage}{\linewidth}\raggedright
% mdparagraph:P269:begin
\subsection*{Cook XII, XIII.35.5}
% mdparagraph:P269:end

% key: KEY-EX-XII-XIII-35-05
% mdparagraph:P270:begin
Hēo ne sang nǣfre bęt ðonne hēo nū dyde.
% mdparagraph:P270:end

% output: EX-XII-XIII-35-05-O001
% mdparagraph:P271:begin
\textbf{Meaning:} She never sang better than she did now.
% mdparagraph:P271:end

% mdparagraph:P272:begin
\textbf{Dyde} repeats the singing action, this time in the third singular. Keep \textbf{nū}, the now of the narrated viewpoint, rather than mechanically changing it to then. The comparative is adverbial \textbf{bęt}. Repeating sang could convey the comparison, but would not display the selected substitute-action pattern.
% mdparagraph:P272:end

\end{minipage}\par\addvspace{4pt}

\par\begin{minipage}{\linewidth}\raggedright
% mdparagraph:P273:begin
\subsection*{Cook XII, XIII.35.6}
% mdparagraph:P273:end

% key: KEY-EX-XII-XIII-35-06
% mdparagraph:P274:begin
Ðrīwa ic ðolode forlidennesse; āne niht and ānne dæg ic wæs on ðissum sǣ.
% mdparagraph:P274:end

% output: EX-XII-XIII-35-06-O001
% mdparagraph:P275:begin
\textbf{Meaning:} Three times I suffered shipwreck; I was in this sea for a night and a day.
% mdparagraph:P275:end

% mdparagraph:P276:begin
\textbf{Ðrīwa} counts the sufferings; \textbf{forlidennesse} is what was suffered. The time extents \textbf{āne niht} and \textbf{ānne dæg} are accusative. \textbf{On ðissum sǣ} is locative dative with the selected masculine noun. Do not make a night and a day two new subjects of wæs.
% mdparagraph:P276:end

\end{minipage}\par\addvspace{4pt}

\par\begin{minipage}{\linewidth}\raggedright
% mdparagraph:P277:begin
\subsection*{Cook XII, XIII.35.7}
% mdparagraph:P277:end

% key: KEY-EX-XII-XIII-35-07
% mdparagraph:P278:begin
Ic wiste sumne mann sē ðe wæs ōð ðone ðriddan heofon lǣded and ðǣr hīerde ðā word ðe nān eorðlic mann ne mōt sprecan.
% mdparagraph:P278:end

% output: EX-XII-XIII-35-07-O001
% mdparagraph:P279:begin
\textbf{Meaning:} I knew a certain man who was carried to the third heaven and heard there words that no earthly man may speak.
% mdparagraph:P279:end

% mdparagraph:P280:begin
\textbf{Wiste} takes the accusative person \textbf{sumne mann}; knowing people is not restricted to cunnan. \textbf{Wæs … lǣded} is a passive carrying event, with the supplied \textbf{ōð} endpoint. The relative hearer is the same man. In the final clause \textbf{nān … ne mōt} expresses prohibition, and \textbf{sprecan} is its plain infinitive. Translating it as a positive permission would reverse the meaning.
% mdparagraph:P280:end

\end{minipage}\par\addvspace{4pt}

\par\begin{minipage}{\linewidth}\raggedright
% mdparagraph:P281:begin
\subsection*{Cook XII, XIII.35.8}
% mdparagraph:P281:end

% key: KEY-EX-XII-XIII-35-08
% mdparagraph:P282:begin
Ēalā, hū nearu is ðæt geat and se weg ðe tō līfe lǣtt, and swīðe fēawe sindon ðe ðone weg findað.
% mdparagraph:P282:end

% output: EX-XII-XIII-35-08-O001
% mdparagraph:P283:begin
\textbf{Meaning:} How narrow is the gate and the way that leads to life, and very few are those who find the way.
% mdparagraph:P283:end

% mdparagraph:P284:begin
The complete Old English \textbf{hū nearu is ðæt geat and se weg} is supplied with its source-shaped singular agreement. It treats the figurative route as a unit and is not a rule for every coordinated subject. \textbf{Lǣtt} stays final in its relative clause. Supplied \textbf{fēawe} is plural and takes \textbf{sindon}; \textbf{swīðe} preserves very, and the requested there is omitted.
% mdparagraph:P284:end

\end{minipage}\par\addvspace{4pt}

\par\begin{minipage}{\linewidth}\raggedright
% mdparagraph:P285:begin
\subsection*{Cook XII, XIII.35.9}
% mdparagraph:P285:end

% key: KEY-EX-XII-XIII-35-09
% mdparagraph:P286:begin
Hīe cumað tō ēow on scēapa gegierelum, ac innan hīe sindon wulfas.
% mdparagraph:P286:end

% output: EX-XII-XIII-35-09-O001
% mdparagraph:P287:begin
\textbf{Meaning:} They come to you in sheep’s garments, but inwardly they are wolves.
% mdparagraph:P287:end

% mdparagraph:P288:begin
\textbf{Cumað} is a present plural statement, not a command. \textbf{On scēapa gegierelum} is the supplied clothing expression: scēapa is genitive plural and gegierelum dative plural. \textbf{Innan} supplies the inward contrast. Treating on as a destination here would change clothes into a place they enter.
% mdparagraph:P288:end

\end{minipage}\par\addvspace{4pt}

\par\begin{minipage}{\linewidth}\raggedright
% mdparagraph:P289:begin
\subsection*{Cook XII, XIII.35.10}
% mdparagraph:P289:end

% key: KEY-EX-XII-XIII-35-10
% mdparagraph:P290:begin
Ðā hē ðǣre byrig geate nēalǣcte, ðā wæs ðǣr dēad mann boren.
% mdparagraph:P290:end

% output: EX-XII-XIII-35-10-O001
% mdparagraph:P291:begin
\textbf{Meaning:} When he approached the city gate, a dead man was being carried there.
% mdparagraph:P291:end

% mdparagraph:P292:begin
The first \textbf{ðā} means when, the second then, and \textbf{ðǣr} means there. \textbf{Nēalǣcte} governs dative \textbf{geate}; \textbf{ðǣre byrig} is its genitive city dependent. \textbf{Wæs … boren} is passive \textbf{beran}, as requested. English being carried is a natural rendering of that event; it does not require an additional Old English progressive verb.
% mdparagraph:P292:end

\end{minipage}\par\addvspace{4pt}

\par\begin{minipage}{\linewidth}\raggedright
% mdparagraph:P293:begin
\subsection*{Cook XII, XIII.35.11}
% mdparagraph:P293:end

% key: KEY-EX-XII-XIII-35-11
% mdparagraph:P294:begin
Ic gesēo mǣden geond ðā strǣte iernan.
% mdparagraph:P294:end

% output: EX-XII-XIII-35-11-O001
% mdparagraph:P295:begin
\textbf{Meaning:} I see a girl running along the street.
% mdparagraph:P295:end

% mdparagraph:P296:begin
Accusative \textbf{mǣden} is the actor of perception infinitive \textbf{iernan}. \textbf{Geond ðā strǣte} takes accusative and identifies the space traversed. The model uses the selected Cook headword mǣden, without an added g, and the accusative strǣte rather than the nominative strǣt.
% mdparagraph:P296:end

\end{minipage}\par\addvspace{4pt}

\par\begin{minipage}{\linewidth}\raggedright
% mdparagraph:P297:begin
\subsection*{Cook XII, XIII.35.12}
% mdparagraph:P297:end

% key: KEY-EX-XII-XIII-35-12
% mdparagraph:P298:begin
Smiere mē mid ęle.
% mdparagraph:P298:end

% output: EX-XII-XIII-35-12-O001
% mdparagraph:P299:begin
\textbf{Meaning:} Anoint me with oil.
% mdparagraph:P299:end

% mdparagraph:P300:begin
\textbf{Smiere} is the selected singular command, with accusative \textbf{mē} and material \textbf{mid ęle}. The spelling \textbf{smire} is independently attested and can be valid Old English; the selected table keeps ie. Do not omit the singular -e merely because other commands such as dēm lack it.
% mdparagraph:P300:end

\end{minipage}\par\addvspace{4pt}

\par\begin{minipage}{\linewidth}\raggedright
% mdparagraph:P301:begin
\subsection*{Cook XII, XIII.35.13}
% mdparagraph:P301:end

% key: KEY-EX-XII-XIII-35-13
% mdparagraph:P302:begin
Sęllað ūs of ēowrum ęle.
% mdparagraph:P302:end

% output: EX-XII-XIII-35-13-O001
% mdparagraph:P303:begin
\textbf{Meaning:} Give us some of your oil.
% mdparagraph:P303:end

% mdparagraph:P304:begin
The selected model addresses several people: \textbf{Sęllað} and \textbf{ēowrum} agree in that reference. \textbf{Ūs} is the dative recipient, and supplied \textbf{of ēowrum ęle} is partitive, not departure from a place. For one addressee, the already taught singular command and possessive give \textbf{Sęle ūs of ðīnum ęle}, using the same partitive help. Keep command and possessive reference consistent. This is an alternative to the selected plural model; the English itself does not spell out singular versus plural you.
% mdparagraph:P304:end

\end{minipage}\par\addvspace{4pt}

\par\begin{minipage}{\linewidth}\raggedright
% mdparagraph:P305:begin
\subsection*{Cook XII, XIII.35.14}
% mdparagraph:P305:end

% key: KEY-EX-XII-XIII-35-14
% mdparagraph:P306:begin
Gif se Sunu ēow ālīest, gē bēoð sōðlice frige.
% mdparagraph:P306:end

% output: EX-XII-XIII-35-14-O001
% mdparagraph:P307:begin
\textbf{Meaning:} If the Son frees you, you will indeed be free.
% mdparagraph:P307:end

% mdparagraph:P308:begin
\textbf{Ālīest} is third singular; \textbf{ēow} is the plural person freed, and \textbf{gē bēoð} preserves that plural reference without shall. The predicate plural \textbf{frige} is supplied and directly supported. The dictionary headword frēo is not the selected plural answer.
% mdparagraph:P308:end

\end{minipage}\par\addvspace{4pt}

\clearpage
% mdparagraph:P309:begin
\section*{Cook XII, XIII.36}
% mdparagraph:P309:end

\par\begin{minipage}{\linewidth}\raggedright
% mdparagraph:P310:begin
\subsection*{Cook XII, XIII.36.1}
% mdparagraph:P310:end

% key: KEY-EX-XII-XIII-36-01
% mdparagraph:P311:begin
Ne dēme gē be ansīene, ac dēmað rihtne dōm.
% mdparagraph:P311:end

% output: EX-XII-XIII-36-01-O001
% mdparagraph:P312:begin
\textbf{Meaning:} Do not judge according to appearance; judge rightly.
% mdparagraph:P312:end

% mdparagraph:P313:begin
Before the following subject \textbf{gē}, the selected plural command is shortened \textbf{dēme}. In the second clause, without a following subject pronoun, the ordinary plural is \textbf{dēmað}. Both remain commands to several people. \textbf{Be ansīene} is the supplied according-to expression; \textbf{rihtne dōm} is the accusative cognate object, “right judgment.” Cook’s positional target is specific; do not claim that a different historical variant could never occur.
% mdparagraph:P313:end

\end{minipage}\par\addvspace{4pt}

\par\begin{minipage}{\linewidth}\raggedright
% mdparagraph:P314:begin
\subsection*{Cook XII, XIII.36.2}
% mdparagraph:P314:end

% key: KEY-EX-XII-XIII-36-02
% mdparagraph:P315:begin
Ðā bewęnde hē him tō ðǣm wīfe and cwæð, Gesiehst ðū ðis wīf? Ic ēode intō ðīnum hūse; ðū ne sealdest mē wæter tō mīnum fōtum; ðēos ðwōh mīne fēt mid hiere tēarum.
% mdparagraph:P315:end

% output: EX-XII-XIII-36-02-O001
% mdparagraph:P316:begin
\textbf{Meaning:} He turned to the woman and said, “Do you see this woman? I entered your house; you gave me no water for my feet, but she washed my feet with her tears.”
% mdparagraph:P316:end

% mdparagraph:P317:begin
\textbf{Bewęnde hē him} has the specifically taught dative reflexive, and \textbf{tō ðǣm wīfe} is the goal. \textbf{Gesiehst} is the selected second-singular form. The supplied feminine \textbf{ðēos} refers to the woman as a person despite the grammatical neuter wīf. \textbf{Ðwōh} is third-singular past; \textbf{mīne fēt} is its accusative object and \textbf{mid hiere tēarum} the washing liquid. Keep the separate supplied intended-use phrase \textbf{wæter tō mīnum fōtum}; it is not another movement destination.
% mdparagraph:P317:end

\end{minipage}\par\addvspace{4pt}

\par\begin{minipage}{\linewidth}\raggedright
% mdparagraph:P318:begin
\subsection*{Cook XII, XIII.36.3}
% mdparagraph:P318:end

% key: KEY-EX-XII-XIII-36-03
% mdparagraph:P319:begin
Ðū mē coss ne sealdest; ðēos, siððan ic inēode, ne geswāc ðæt hēo mīne fēt ne cyste.
% mdparagraph:P319:end

% output: EX-XII-XIII-36-03-O001
% mdparagraph:P320:begin
\textbf{Meaning:} You gave me no kiss; since I entered, this woman has not stopped kissing my feet.
% mdparagraph:P320:end

% mdparagraph:P321:begin
The first clause ends in \textbf{sealdest}, as directed, and contains no indefinite article. \textbf{Inēode} is the selected first-singular past of ingān. The entire \textbf{ne geswāc ðæt hēo mīne fēt ne cyste} is supplied from the attested negative cessation construction. Both negatives belong to it: the inner ne does not mean she refused to kiss. Do not add a have auxiliary or invent a general rule for every complement of cease.
% mdparagraph:P321:end

\end{minipage}\par\addvspace{4pt}

\par\begin{minipage}{\linewidth}\raggedright
% mdparagraph:P322:begin
\subsection*{Cook XII, XIII.36.4}
% mdparagraph:P322:end

% key: KEY-EX-XII-XIII-36-04
% mdparagraph:P323:begin
Mīn hēafod ðū mid ęlum ne smieredest; ðēos smierede mīne fēt mid sealfe.
% mdparagraph:P323:end

% output: EX-XII-XIII-36-04-O001
% mdparagraph:P324:begin
\textbf{Meaning:} You did not anoint my head with oils; this woman anointed my feet with ointment.
% mdparagraph:P324:end

% mdparagraph:P325:begin
Keep the instructed plural \textbf{ęlum} and singular \textbf{sealfe}. Second-singular \textbf{smieredest} has -st; third-singular \textbf{smierede} does not. The objects hēafod and fēt are accusative, and both materials follow mid in the dative. Attested i/\allowbreak{}y spellings can be linguistic alternatives, but cannot remove the person distinction.
% mdparagraph:P325:end

\end{minipage}\par\addvspace{4pt}

\par\begin{minipage}{\linewidth}\raggedright
% mdparagraph:P326:begin
\subsection*{Cook XII, XIII.36.5}
% mdparagraph:P326:end

% key: KEY-EX-XII-XIII-36-05
% mdparagraph:P327:begin
For-ðon ic ðē sęcge, Maniga synna hiere sindon forgiefen, for-ðon-ðe hēo mē swīðe lufode.
% mdparagraph:P327:end

% output: EX-XII-XIII-36-05-O001
% mdparagraph:P328:begin
\textbf{Meaning:} Therefore I tell you that many sins are forgiven her, because she loved me greatly.
% mdparagraph:P328:end

% mdparagraph:P329:begin
\textbf{Synna} is the plural passive subject, \textbf{hiere} the dative person forgiven, and \textbf{forgiefen} an uninflected predicative participle. It is not a finite plural verb. The opening supplied \textbf{For-ðon} means therefore; \textbf{for-ðon-ðe} introduces the reason. Keep \textbf{lufode} last in that clause and retain swīðe. Merely knowing forgiefan as give/\allowbreak{}grant would not establish this forgiveness sense.
% mdparagraph:P329:end

\end{minipage}\par\addvspace{4pt}

\par\begin{minipage}{\linewidth}\raggedright
% mdparagraph:P330:begin
\subsection*{Cook XII, XIII.36.6}
% mdparagraph:P330:end

% key: KEY-EX-XII-XIII-36-06
% mdparagraph:P331:begin
Lǣs lufað ðǣm ðe lǣs forgiefen is.
% mdparagraph:P331:end

% output: EX-XII-XIII-36-06-O001
% mdparagraph:P332:begin
\textbf{Meaning:} The person to whom less is forgiven loves less.
% mdparagraph:P332:end

% mdparagraph:P333:begin
The complete compressed relative sentence is supplied. The first \textbf{lǣs} is the prompt’s neuter accusative degree expression, “less” with loves. The second \textbf{lǣs} is nominative neuter singular: the smaller amount that is forgiven. \textbf{Ðǣm ðe} is dative within the forgiveness relation, with the matrix subject understood. This is not evidence that lufian normally takes a dative subject. An expanded relative could express the meaning, but would replace the supplied compact target.
% mdparagraph:P333:end

\end{minipage}\par\addvspace{4pt}

\par\begin{minipage}{\linewidth}\raggedright
% mdparagraph:P334:begin
\subsection*{Cook XII, XIII.36.7}
% mdparagraph:P334:end

% key: KEY-EX-XII-XIII-36-07
% mdparagraph:P335:begin
Hit ne ālīefð ðē hīe tō wīfe habban.
% mdparagraph:P335:end

% output: EX-XII-XIII-36-07-O001
% mdparagraph:P336:begin
\textbf{Meaning:} You are not permitted to have her as your wife.
% mdparagraph:P336:end

% mdparagraph:P337:begin
The whole permission frame is supplied: impersonal \textbf{hit ne ālīefð}, dative permitted person \textbf{ðē}, and \textbf{hīe} as accusative object of \textbf{habban}. \textbf{Tō wīfe} expresses her role as wife, not a destination. The English is is omitted because the finite permitting verb already supplies the predicate; do not add wesan.
% mdparagraph:P337:end

\end{minipage}\par\addvspace{4pt}

\par\begin{minipage}{\linewidth}\raggedright
% mdparagraph:P338:begin
\subsection*{Cook XII, XIII.36.8}
% mdparagraph:P338:end

% key: KEY-EX-XII-XIII-36-08
% mdparagraph:P339:begin
Ðā gehēt hē mid āðe hiere ðæt tō sęllanne, swā hwæs swā hēo hine bǣde.
% mdparagraph:P339:end

% output: EX-XII-XIII-36-08-O001
% mdparagraph:P340:begin
\textbf{Meaning:} He then promised with an oath to give her whatever she asked him for.
% mdparagraph:P340:end

% mdparagraph:P341:begin
Use the supplied promise/\allowbreak{}free-choice bridge. \textbf{Hiere} is the dative giving recipient and resumptive \textbf{ðæt} the accusative thing given. Inside \textbf{swā hwæs swā}, biddan requires genitive \textbf{hwæs} for the thing requested and accusative \textbf{hine} for the person asked. \textbf{Bǣde} is preterite subjunctive; the omitted English should does not become another verb. One English whatever cannot force the same case into two different constructions.
% mdparagraph:P341:end

\end{minipage}\par\addvspace{4pt}

\par\begin{minipage}{\linewidth}\raggedright
% mdparagraph:P342:begin
\subsection*{Cook XII, XIII.36.9}
% mdparagraph:P342:end

% key: KEY-EX-XII-XIII-36-09
% mdparagraph:P343:begin
Ðā cwæð hēo, fram hiere mēder myngod, Sęle mē on disce Johannes ðæs fuluhteres hēafod.
% mdparagraph:P343:end

% output: EX-XII-XIII-36-09-O001
% mdparagraph:P344:begin
\textbf{Meaning:} Prompted by her mother, she said, “Give me John the Baptist’s head on a dish.”
% mdparagraph:P344:end

% mdparagraph:P345:begin
Supplied \textbf{myngod} is an uninflected predicative participle; \textbf{fram hiere mēder} supplies its agent. Singular \textbf{Sęle} has dative \textbf{mē} and accusative \textbf{hēafod}. Supplied \textbf{Johannes} is genitive here, with matching \textbf{ðæs fuluhteres}. The primary name spelling has initial I; the declared J follows Cook’s headword. That spelling choice is separate from case.
% mdparagraph:P345:end

\end{minipage}\par\addvspace{4pt}

\par\begin{minipage}{\linewidth}\raggedright
% mdparagraph:P346:begin
\subsection*{Cook XII, XIII.36.10}
% mdparagraph:P346:end

% key: KEY-EX-XII-XIII-36-10
% mdparagraph:P347:begin
Ðā āsęnde hē and behēafdode Johannem on ðǣm cwearterne.
% mdparagraph:P347:end

% output: EX-XII-XIII-36-10-O001
% mdparagraph:P348:begin
\textbf{Meaning:} He then sent and beheaded John in the prison.
% mdparagraph:P348:end

% mdparagraph:P349:begin
The source leaves the sending patient unexpressed; the model preserves that bounded absolute \textbf{āsęnde}. \textbf{Behēafdode} is third-singular past and takes supplied accusative \textbf{Johannem}. The prison is locative \textbf{on ðǣm cwearterne}. Do not keep the name in its dictionary nominative merely because English John shows no case ending, or silently replace the active beheaded with a different command construction.
% mdparagraph:P349:end

\end{minipage}\par\addvspace{4pt}

\par\begin{minipage}{\linewidth}\raggedright
% mdparagraph:P350:begin
\subsection*{Cook XII, XIII.36.11}
% mdparagraph:P350:end

% key: KEY-EX-XII-XIII-36-11
% mdparagraph:P351:begin
And man brōhte his hēafod on disce.
% mdparagraph:P351:end

% output: EX-XII-XIII-36-11-O001
% mdparagraph:P352:begin
\textbf{Meaning:} Someone brought his head on a dish.
% mdparagraph:P352:end

% mdparagraph:P353:begin
Indefinite \textbf{man} supplies the unnamed human agent, and \textbf{brōhte} is third-singular past. \textbf{His hēafod} is accusative object; \textbf{on disce} identifies the dish carrying it. English one means someone here, not a numeral modifying the head. This is a separate final narrated event and must not be omitted.
% mdparagraph:P353:end

\end{minipage}\par\addvspace{4pt}

\clearpage
% mdparagraph:P354:begin
\section*{R. The gymnasium scene}
% mdparagraph:P354:end

\reserve{8\baselineskip}
% mdparagraph:P355:begin
\subsection*{R01}
% mdparagraph:P355:end

% key: KEY-NEW-U16-R01
% output: NEW-U16-R01-O001
% mdparagraph:P356:begin
\textbf{1. Whole translation.} After these words he went along the way that had been shown to him, until he reached the city gate and entered there. While he was thinking whom he might ask for help in sustaining his life, he saw a naked boy running along the street. The boy was anointed with oil and girded with a cloth, and carried in his hand the equipment for young men’s play belonging to the baths. He called out in a loud voice and said: “Listen, citizens! Listen, strangers, free and enslaved, noble and lowborn! The baths are open.”
% mdparagraph:P356:end

% mdparagraph:P357:begin
When Apollonius heard this, he took off the half-cloak that he was wearing and went into the bathing place. As he observed each of them at their activity, he looked for his equal, but could not find him there in the crowd. Then suddenly Arcestrates, king of that whole country, arrived with a large company of his men and entered the bath. The king began playing ball with his companions. Apollonius joined in the king’s game, as God willed; running, he caught the ball and, after striking it with swift speed, sent it back to the king who was playing. The king sent it back again; Apollonius struck it swiftly, so that he never let it fall. The king then recognized the young man’s agility, so that he knew that he had no equal in that game. He then said to his companions: “Go away from here; this young man, as it seems to me, is my equal.”
% mdparagraph:P357:end

% mdparagraph:P358:begin
When Apollonius heard that the king was praising him, he ran quickly and approached the king, and with a skilled hand he swung the \emph{top}—apparently the ball—with such great speed that it seemed to the king as though he had been changed from old age back to youth. Afterwards Apollonius attended the king agreeably at his royal seat. When the king came out of the bath, Apollonius led him by the hand, and then Apollonius went away from there along the way by which he had come earlier.
% mdparagraph:P358:end

% mdparagraph:P359:begin
After Apollonius had gone away, the king said to his men: “I swear by our common safety that I have never bathed better than I did today, through the service of a young man whose identity I do not know.” He then turned his gaze towards one of his men and said: “Go and find out what sort of man the young man is who served me so well today.”
% mdparagraph:P359:end

% mdparagraph:P360:begin
The man then went after Apollonius. When he saw that Apollonius was wrapped in a dirty cloak, he returned to the king and said: “The young man you asked about is a shipwrecked man.” The king said: “How do you know that?” The man answered him: “Even if he himself keeps silent about it, his clothing reveals his condition.” Then the king said: “Go quickly and tell him that ‘the king asks you to come to his meal.’”
% mdparagraph:P360:end

% mdparagraph:P361:begin
This model retains every stage from arrival to invitation. \textbf{Plegan} at 168.6 is grammatical accusative singular; “equipment for young men’s play” is a contextual rendering, not a claim that the Old English noun is plural. “Bath,” “baths” or “bathing place” can render the facility without changing the action.
% mdparagraph:P361:end

% mdparagraph:P362:begin
The pronouns in the ball exchange are interpreted from the game sequence. At 169.5–7, however, the model deliberately leaves “he knew that he had no equal” unresolved; \textbf{næfde} must stay negative. At 169.11–13, \textbf{top} is retained with Cook’s tentative ball interpretation. “Our common safety,” “common welfare” or a carefully explained religious rendering may express the oath at 170.2; none should add an event. See R06 for the two substantial textual uncertainties.
% mdparagraph:P362:end

% mdparagraph:P363:begin
The final invitation is intentionally awkward in English because the source combines an outer report with quoted second-person address. A fluent paraphrase may explain the intended message, but a complete translation should acknowledge the actual mixture rather than conceal it. The following feast is not part of this reading.
% mdparagraph:P363:end

\reserve{8\baselineskip}
% mdparagraph:P364:begin
\subsection*{R02}
% mdparagraph:P364:end

% key: KEY-NEW-U16-R02
% output: NEW-U16-R02-O001
% mdparagraph:P365:begin
\textbf{1. Entry into the bath.} The boy announces that the baths are open, and Apollonius hears him: \textbf{Se bæðstęde is open} (168.9–10), followed by Apollonius taking off his half-cloak and entering (168.11–12). His earlier concern is finding help to live on; the king has not yet invited him.
% mdparagraph:P365:end

% output: NEW-U16-R02-O002
% mdparagraph:P366:begin
\textbf{2. The game.} Apollonius joins the ball game, catches and returns the ball with speed, and does not let it fall (168.18–169.5). The king notices \textbf{ðæs geongan snelnesse}, the young man’s agility. Merely saying that they played omits the evidence of skill that attracts attention.
% mdparagraph:P366:end

% output: NEW-U16-R02-O003
% mdparagraph:P367:begin
\textbf{3. The king's statement.} The king tells his companions to go away and says \textbf{ðes cniht … is mīn gelīca}, “this young man … is my equal” (169.7–8). Distinguish this clear direct statement from the preceding clause whose pronoun reference Cook calls unclear.
% mdparagraph:P367:end

% output: NEW-U16-R02-O004
% mdparagraph:P368:begin
\textbf{4. The messenger's conclusion.} The messenger concludes that Apollonius is \textbf{forliden mann}, a shipwrecked man (170.10–11). He has seen him wrapped in a dirty cloak; he says \textbf{his gegierela hine gesweotolað}, “his clothing reveals him” (170.13–14). He infers the condition from clothing, rather than reporting a confession or learning the man’s name.
% mdparagraph:P368:end

% output: NEW-U16-R02-O005
% mdparagraph:P369:begin
\textbf{5. Final order.} The king orders his messenger to go quickly and convey an invitation to Apollonius to come to the king’s meal (170.14–16). The messenger is to deliver the request; the text does not yet narrate Apollonius accepting or attending the feast.
% mdparagraph:P369:end

\reserve{8\baselineskip}
% mdparagraph:P370:begin
\subsection*{R03}
% mdparagraph:P370:end

% key: KEY-NEW-U16-R03
% output: NEW-U16-R03-O001
% mdparagraph:P371:begin
\textbf{1. 169.4 hine.} \textbf{Hine} is accusative masculine singular and refers to the \textbf{ball}, ðoðor; Cook’s note 3 confirms this. It is the actor of \textbf{feallan} in the construction “he never let it fall.” Grammatical masculine does not make its referent a man.
% mdparagraph:P371:end

% output: NEW-U16-R03-O002
% mdparagraph:P372:begin
\textbf{2. 169.14 him.} \textbf{Him} is dative singular and refers to the \textbf{king}, the person Apollonius serves with \textbf{ðēnode}. The verb’s recipient/\allowbreak{}service frame and the scene establish its role; it is not an accusative object merely because English says served him.
% mdparagraph:P372:end

% output: NEW-U16-R03-O003
% mdparagraph:P373:begin
\textbf{3. 169.15 hē before hine.} The \textbf{hē} immediately before hine lǣdde is nominative singular, \textbf{Apollonius}, subject of the leading action. Cook’s note 14 identifies him explicitly. It is distinct from the king who has just come out of the bath.
% mdparagraph:P373:end

% output: NEW-U16-R03-O004
% mdparagraph:P374:begin
\textbf{4. 169.15 hine.} \textbf{Hine} is accusative masculine singular, the \textbf{king}, direct object of lǣdde. Cook’s note 15 names the king. Apollonius leads the king by the hand, not the reverse.
% mdparagraph:P374:end

% output: NEW-U16-R03-O005
% mdparagraph:P375:begin
\textbf{5. 169.16 him.} \textbf{Him} is a dative reflexive referring to \textbf{Apollonius}, the subject who then turns away. It belongs to \textbf{gewęnde}, as the local help and §184a indicate. It does not introduce another person to whom Apollonius went.
% mdparagraph:P375:end

% output: NEW-U16-R03-O006
% mdparagraph:P376:begin
\textbf{6. 170.9 hē.} The \textbf{hē} after ðæt is nominative masculine singular, \textbf{Apollonius}, subject of \textbf{wæs … bewǣfed}. The preceding observer is the messenger. Cook’s note 8 and the clothing report distinguish the two referents; the masculine ending alone would not.
% mdparagraph:P376:end

\reserve{8\baselineskip}
% mdparagraph:P377:begin
\subsection*{R04}
% mdparagraph:P377:end

% key: KEY-NEW-U16-R04
% output: NEW-U16-R04-O001
% mdparagraph:P378:begin
\textbf{1. 168.1–2 relation.} \textbf{Ðe him getǣht wæs} is a relative clause identifying \textbf{ðone weg}, the road that had been shown to him. It is not the main action of walking.
% mdparagraph:P378:end

% output: NEW-U16-R04-O002
% mdparagraph:P379:begin
\textbf{2. 168.1–2 finite group.} The passive group is \textbf{getǣht wæs}: \textbf{wæs}, third-singular preterite indicative of wesan, is finite; \textbf{getǣht} is the past participle. The relative’s subject is the road, and \textbf{him} is the dative person to whom it was shown.
% mdparagraph:P379:end

% output: NEW-U16-R04-O003
% mdparagraph:P380:begin
\textbf{3. 168.2 relation.} \textbf{Ōð-ðæt hē becōm} is a temporal endpoint clause: the walking continues until he reaches the gate.
% mdparagraph:P380:end

% output: NEW-U16-R04-O004
% mdparagraph:P381:begin
\textbf{4. 168.2 finite form.} \textbf{Becōm} is third-singular preterite indicative of becuman. The completed arrival supplies the endpoint; do not call it subjunctive merely because it is subordinate.
% mdparagraph:P381:end

% output: NEW-U16-R04-O005
% mdparagraph:P382:begin
\textbf{5. 168.3 relation.} \textbf{Hwone hē biddan meahte} is a dependent question giving the content of his thought: whom he might ask. \textbf{Hwone} denotes the person asked, not the subject of could.
% mdparagraph:P382:end

% output: NEW-U16-R04-O006
% mdparagraph:P383:begin
\textbf{6. 168.3 finite group.} The group is \textbf{biddan meahte}. \textbf{Meahte} is third-singular preterite of magan, selected here as subjunctive possibility in the dependent question; the indicative has the same surface. \textbf{Biddan} is a plain infinitive, not a second finite verb.
% mdparagraph:P383:end

% output: NEW-U16-R04-O007
% mdparagraph:P384:begin
\textbf{7. 169.9 relation.} \textbf{Ðæt se cyning hine hęrede} is the content of what Apollonius heard: that the king praised him. The surrounding when-clause contains the hearing.
% mdparagraph:P384:end

% output: NEW-U16-R04-O008
% mdparagraph:P385:begin
\textbf{8. 169.9 finite form.} \textbf{Hęrede} is third-singular preterite of hęrian, selected here as indicative in the heard content. Its subjunctive spelling is identical. The answer should identify the king as its subject and avoid pretending that the ending by itself proves mood.
% mdparagraph:P385:end

% output: NEW-U16-R04-O009
% mdparagraph:P386:begin
\textbf{9. 170.6 relation.} \textbf{Hwæt se geonga mann sīe} is the dependent question governed by \textbf{gewite}, “find out.” Here hwæt asks what the man is or what his condition is; the messenger responds with his shipwrecked condition.
% mdparagraph:P386:end

% output: NEW-U16-R04-O010
% mdparagraph:P387:begin
\textbf{10. 170.6 finite form.} \textbf{Sīe} is third-singular present subjunctive of wesan. The subject is \textbf{se geonga mann}. It is finite, whereas gewite is the separate imperative governing the question.
% mdparagraph:P387:end

% output: NEW-U16-R04-O011
% mdparagraph:P388:begin
\textbf{11. 170.13 relation.} \textbf{Ðēah hē hit self forswīge} is a concessive clause: even if he himself keeps it silent. It concedes possible silence while asserting that the clothing still reveals him.
% mdparagraph:P388:end

% output: NEW-U16-R04-O012
% mdparagraph:P389:begin
\textbf{12. 170.13 finite form.} \textbf{Forswīge} is third-singular present subjunctive of forswīgian in the supplied source construction. \textbf{Hē} is the subject and \textbf{hit} the concealed matter; \textbf{self} reinforces the subject rather than supplying another finite verb.
% mdparagraph:P389:end

\reserve{8\baselineskip}
% mdparagraph:P390:begin
\subsection*{R05}
% mdparagraph:P390:end

% key: KEY-NEW-U16-R05
% output: NEW-U16-R05-O001
% mdparagraph:P391:begin
\textbf{1. Outer speaker.} The \textbf{king, Arcestrates}, is the outer speaker, introduced by \textbf{Ðā cwæð se cyning}.
% mdparagraph:P391:end

% output: NEW-U16-R05-O002
% mdparagraph:P392:begin
\textbf{2. Immediate addressee.} His \textbf{messenger}, the man who has just returned with the report, is the immediate addressee of singular \textbf{Gā} and \textbf{sęge}. Apollonius is elsewhere at this point.
% mdparagraph:P392:end

% output: NEW-U16-R05-O003
% mdparagraph:P393:begin
\textbf{3. Him.} \textbf{Him} denotes \textbf{Apollonius}, the dative recipient of the message after sęge.
% mdparagraph:P393:end

% output: NEW-U16-R05-O004
% mdparagraph:P394:begin
\textbf{4. Se cyning.} \textbf{Se cyning} inside the proposed message denotes \textbf{the same king, Arcestrates}. The king gives the messenger wording that names him in the third person.
% mdparagraph:P394:end

% output: NEW-U16-R05-O005
% mdparagraph:P395:begin
\textbf{5. Ðē and ðū.} Inner \textbf{ðē} and \textbf{ðū} both denote \textbf{Apollonius}, the person who will hear the invitation. \textbf{Ðē} is the accusative person whom the king asks; \textbf{ðū} is nominative subject of \textbf{cume}. Their referent is not the messenger merely because the king is currently addressing him.
% mdparagraph:P395:end

% output: NEW-U16-R05-O006
% mdparagraph:P396:begin
\textbf{6. Construction.} The outer \textbf{sęge him ðæt} begins a report, but what follows is quoted wording with second-person address: \textbf{se cyning bitt ðē ðæt ðū cume tō his gereorde}. Cook’s note 15 calls this a mixture of direct and indirect construction. Preserve that mixture when explaining the source; a smooth indirect paraphrase can explain its meaning but cannot replace the wording being analysed. \textbf{Cume} is present subjunctive in the request; English can render the request with “asks you to come.”
% mdparagraph:P396:end

\reserve{8\baselineskip}
% mdparagraph:P397:begin
\subsection*{R06}
% mdparagraph:P397:end

% key: KEY-NEW-U16-R06
% output: NEW-U16-R06-O001
% mdparagraph:P398:begin
\textbf{1. 169.5–7 translation.} “The king then recognized the young man’s agility, so that he knew that he had no equal in that game.” This preserves the repeated pronouns and the negative \textbf{næfde} rather than silently deciding all their referents.
% mdparagraph:P398:end

% output: NEW-U16-R06-O002
% mdparagraph:P399:begin
\textbf{2. 169.5–7 limit.} Original note 6 says the clause is unclear. Its Latin parallel has the king knowing that he himself had no equal; the Old English’s sequence and repeated hē/\allowbreak{}his leave room for discussion, including a reading in which the young man is the one without an equal. A proposed reading should name its referents and acknowledge that limit. Dropping ne from næfde, or changing the whole clause to “he had found his equal,” is not a translation of the printed wording.
% mdparagraph:P399:end

% output: NEW-U16-R06-O003
% mdparagraph:P400:begin
\textbf{3. 169.11–13 translation.} “With a skilled hand he swung the top—apparently the ball—with such great speed that it seemed to the king as though he had been changed from old age back to youth.” The king experiences the seeming return to youth; Apollonius performs the action.
% mdparagraph:P400:end

% output: NEW-U16-R06-O004
% mdparagraph:P401:begin
\textbf{4. 169.11–13 limit.} Original note 10 says that the Old English departs from the Latin. The English help for the Latin describes rubbing/\allowbreak{}anointing the king with ointment so rapidly that he seemed young again. Cook only tentatively identifies \textbf{top} with \textbf{ðoðor}. Translate \textbf{swang ðone top}, mark the uncertain object, and do not substitute the ointment action for the Old English event.
% mdparagraph:P401:end

\end{document}
